\documentclass[11pt,a4paper]{article}

\usepackage[T1]{fontenc}
\usepackage[utf8]{inputenc}
\usepackage[polish]{babel}
\usepackage{lmodern}
\usepackage{microtype}
\usepackage{amsmath,amssymb,mathtools}
\usepackage{geometry}
\usepackage{enumitem}
\usepackage{hyperref}

\geometry{margin=2.5cm}
\setlength{\parindent}{0pt}
\setlength{\parskip}{0.45em}
\setcounter{secnumdepth}{0}
\setcounter{tocdepth}{2}
\setlist[itemize]{leftmargin=*,itemsep=0.2em,topsep=0.3em}
\hypersetup{colorlinks=true,linkcolor=blue,urlcolor=blue}

\newcommand{\R}{\mathbb{R}}
\newcommand{\N}{\mathbb{N}}
\newcommand{\dd}{\,\mathrm{d}}

\title{Opracowanie tematów na egzamin licencjacki z fizyki}
\author{}
\date{}

\begin{document}
\maketitle
\tableofcontents
\clearpage

\section{Fizyka ogólna}
\subsection{1. Zasady dynamiki Newtona i zasada zachowania pędu}
\paragraph{Zakres.} Pierwsza, druga i trzecia zasada dynamiki Newtona; pęd, popęd siły i zachowanie pędu w układzie izolowanym.

Pierwsza zasada Newtona definiuje układ inercjalny: jeżeli wypadkowa siła działająca na ciało jest zerowa, ciało pozostaje w spoczynku albo porusza się ruchem jednostajnym prostoliniowym. Druga zasada wiąże zmianę ruchu z oddziaływaniem. Dla stałej masy zapisujemy ją jako
\[
\sum_i \mathbf F_i=\frac{d\mathbf p}{dt}=m\mathbf a,\qquad \mathbf p=m\mathbf v,
\]
gdzie $\mathbf F_i$ są siłami zewnętrznymi, $\mathbf p$ jest pędem, $m$ masą, a $\mathbf a$ przyspieszeniem. Ogólniejsza postać z pochodną pędu jest potrzebna, gdy masa układu się zmienia; dla ciała o zmiennej masie trzeba przy tym jawnie uwzględnić strumień pędu materii. Trzecia zasada mówi, że siły wzajemnego oddziaływania dwóch ciał są równe co do wartości, przeciwnie skierowane i leżą na tej samej prostej. Działają na różne ciała, więc nie znoszą się w równaniu ruchu jednego z nich.

Całka siły po czasie to popęd, równy zmianie pędu: $\mathbf J=\int_{t_1}^{t_2}\mathbf F\,dt=\Delta\mathbf p$. Dla układu wielu ciał siły wewnętrzne parami się znoszą, o ile obowiązuje mechanika Newtona, zatem $d\mathbf P/dt=\mathbf F_{\rm zewn}$, gdzie $\mathbf P=\sum_j\mathbf p_j$. Przy zerowej wypadkowej sile zewnętrznej pęd całkowity jest stały. Przykładem jest zderzenie dwóch wózków na poziomej, mało tarciowej szynie: przed i po zderzeniu suma $m_1\mathbf v_1+m_2\mathbf v_2$ jest taka sama. Zachowanie pędu nie wymaga zachowania energii kinetycznej: zderzenie niesprężyste może tę energię zamienić w ciepło i odkształcenia.

W zadaniu egzaminacyjnym warto najpierw wybrać badany układ i narysować siły zewnętrzne, zamiast od razu zakładać zachowanie pędu. Jeśli impuls siły zewnętrznej w czasie zderzenia jest mały względem pędów ciał, pęd można uznać za zachowany w przybliżeniu nawet wtedy, gdy później pojawi się tarcie. W jednowymiarowym zderzeniu kierunki koduje się znakami prędkości. Trzecia zasada pomaga wyjaśnić wymianę pędu między ciałami, ale sama nie gwarantuje zachowania pędu całego układu, jeśli oddziałuje on z otoczeniem.

\subsection{2. Układy inercjalne i nieinercjalne, siły pozorne}
\paragraph{Zakres.} Układy odniesienia, transformacja Galileusza oraz opis ruchu w układzie przyspieszonym i obracającym się.

Układ inercjalny to taki, w którym ciało swobodne nie przyspiesza. W mechanice klasycznej wszystkie układy poruszające się względem siebie ruchem jednostajnym prostoliniowym są inercjalne. Jeśli układ $S'$ przesuwa się względem $S$ ze stałą prędkością $\mathbf V$, to transformacja Galileusza ma postać $\mathbf r'=\mathbf r-\mathbf Vt$, $t'=t$; przyspieszenie jest takie samo, więc zasady Newtona zachowują formę. Układ związany z przyspieszającym pojazdem nie jest inercjalny. Aby w nim nadal stosować równanie Newtona, dodaje się siłę pozorną przeciwną do przyspieszenia układu: $\mathbf F_{\rm poz}=-m\mathbf A$, gdzie $\mathbf A$ jest przyspieszeniem początku układu.

W układzie obracającym się z prędkością kątową $\boldsymbol\Omega$ pojawiają się dodatkowe składniki. Siła odśrodkowa ma postać $-m\boldsymbol\Omega\times(\boldsymbol\Omega\times\mathbf r)$, a siła Coriolisa dla prędkości $\mathbf v'$ względem układu wynosi $-2m\boldsymbol\Omega\times\mathbf v'$. Przy zmiennej prędkości obrotu występuje także składnik Eulera $-m\dot{\boldsymbol\Omega}\times\mathbf r$. Znak i kierunek łatwo sprawdzić w układzie obracającym się równomiernie: odśrodkowy człon jest skierowany od osi. Na Ziemi siła Coriolisa odchyla ruchy poziome i ma znaczenie dla atmosfery oraz oceanów; w małych, krótkotrwałych doświadczeniach jej wpływ często pomijalny. Siły pozorne nie odpowiadają oddziaływaniu z innym ciałem i zależą od wybranego układu; ich wprowadzenie pozwala opisać obserwacje nieinercjalne, ale nie zmienia fizycznego oddziaływania.

Praktyczny test polega na sprawdzeniu, czy układ współrzędnych porusza się ruchem jednostajnym i prostoliniowym względem układu, w którym obowiązuje zasada bezwładności. W windzie przyspieszającej ku górze pozorny ciężar pasażera rośnie, a przy przyspieszeniu ku dołowi maleje; odczucie nacisku jest związane z siłą podłoża, nie z nową siłą grawitacyjną. W rozwiązaniu należy podać, względem którego układu definiuje się prędkość i przyspieszenie. Pozwala to uniknąć mieszania rzeczywistych sił oddziaływania z członami dodanymi wyłącznie dla wygody opisu.

\subsection{3. Praca, energia kinetyczna i potencjalna, siły zachowawcze}
\paragraph{Zakres.} Praca siły, twierdzenie o energii kinetycznej, energia potencjalna i zachowanie energii mechanicznej.

Praca siły $\mathbf F$ przy przesunięciu po torze $C$ jest całką liniową $W=\int_C\mathbf F\cdot d\mathbf r$. Dla stałej siły i prostoliniowego przesunięcia $W=Fs\cos\theta$, gdzie $\theta$ jest kątem między siłą i przemieszczeniem. Moc chwilowa wynosi $P=dW/dt=\mathbf F\cdot\mathbf v$. Z drugiej zasady Newtona wynika twierdzenie o pracy i energii: praca wypadkowej siły jest zmianą energii kinetycznej, $W_{\rm wypadk}= \Delta K$, przy czym $K=\tfrac12mv^2$ dla nierelatywistycznej cząstki o stałej masie.

Siła jest zachowawcza, jeżeli jej praca między dwoma punktami nie zależy od drogi, równoważnie praca po drodze zamkniętej jest zerowa. W obszarze bez osobliwości można ją zapisać jako $\mathbf F=-\nabla U$, gdzie $U$ jest energią potencjalną określoną z dokładnością do stałej. Dla ciężaru blisko powierzchni Ziemi $U=mgh$; dla sprężyny idealnej $U=\tfrac12kx^2$. Gdy działają wyłącznie siły zachowawcze, $K+U=E_{\rm mech}$ pozostaje stałe: spadająca piłka traci energię potencjalną i zyskuje kinetyczną. Tarcie jest zwykle niezachowawcze; praca tarcia zmniejsza energię mechaniczną i przekształca ją w energię wewnętrzną. Bilans całkowitej energii, obejmujący ciepło i odkształcenia, nadal obowiązuje. W zadaniach trzeba wskazać układ, poziom odniesienia dla $U$ oraz założenia, na przykład zaniedbanie oporu powietrza.

Dobrym sposobem kontroli rachunku jest sprawdzenie jednostek: praca i obie energie mają jednostkę dżula, a moc wat. Zmiana stałej dodanej do energii potencjalnej nie wpływa na siłę ani na obserwowalną różnicę energii, dlatego poziom zerowy wybiera się dogodnie. Jeśli siła zależy od położenia, pracę oblicza się przez całkowanie jej składowej wzdłuż drogi, a nie przez proste mnożenie wartości siły i całego przemieszczenia. Przy obecności tarcia można nadal stosować bilans energii, o ile uwzględni się energię cieplną.

\subsection{4. Dynamika bryły sztywnej i zachowanie momentu pędu}
\paragraph{Zakres.} Ruch obrotowy bryły, moment siły, moment bezwładności, energia obrotu oraz bilans momentu pędu.

Bryła sztywna zachowuje odległości między swoimi punktami. W ruchu obrotowym wokół ustalonej osi kątowe odpowiedniki wielkości liniowych to kąt $\theta$, prędkość kątowa $\omega=d\theta/dt$ i przyspieszenie kątowe $\alpha=d\omega/dt$. Moment siły względem punktu $O$ wynosi $\boldsymbol\tau=\mathbf r\times\mathbf F$. Dla stałej osi i stałego momentu bezwładności $I$ równanie ruchu jest $\tau_{\rm oś}=I\alpha$, analogiczne do $F=ma$. Moment bezwładności $I=\sum_jm_jr_{\perp j}^2$ albo $I=\int r_\perp^2\,dm$ zależy od rozkładu masy względem osi. Energia obrotu to $K_{\rm obr}=\tfrac12I\omega^2$.

Moment pędu punktu materialnego to $\mathbf L=\mathbf r\times\mathbf p$, a bryły obracającej się wokół osi symetrii $L=I\omega$. Ogólnie bilans względem ustalonego inercjalnego początku ma postać $d\mathbf L/dt=\boldsymbol\tau_{\rm zewn}$. Jeżeli zewnętrzny moment siły jest zerowy, moment pędu jest zachowany. Łyżwiarka przyciągająca ramiona zmniejsza $I$, więc jej $\omega$ rośnie, aby zachować $L$; energia obrotowa może wzrosnąć dzięki wykonanej przez nią pracy. Toczenie bez poślizgu spełnia $v_{\rm CM}=R\omega$ i łączy ruch postępowy z obrotowym, lecz warunek ten nie wynika z samej dynamiki, tylko z kinematycznego braku poślizgu. Równanie $\tau=I\alpha$ w tej postaci stosuje się dla ustalonej osi i stałego $I$; w ogólnym ruchu przestrzennym należy korzystać z wektorowego momentu pędu i tensora bezwładności.

Moment siły zależy od wybranego punktu odniesienia, dlatego w obliczeniach trzeba go konsekwentnie wskazać. Siła przyłożona blisko osi daje mniejszy moment niż taka sama siła przyłożona na długim ramieniu, co wykorzystuje się na przykład przy odkręcaniu śruby kluczem. W ruchu obrotowym praca elementarna jest równa momentowi razy przyrostowi kąta, a moc obrotowa jest iloczynem momentu i prędkości kątowej. Dla bryły złożonej wygodnie rozkłada się masę na elementy i sumuje ich wkłady do momentu bezwładności.

\subsection{5. Oscylator harmoniczny, tłumienie, drgania swobodne i wymuszone}
\paragraph{Zakres.} Ruch harmoniczny prosty, energia oscylatora, tłumienie oraz odpowiedź na wymuszenie i rezonans.

Oscylator harmoniczny powstaje, gdy siła powrotna jest proporcjonalna do wychylenia i skierowana ku położeniu równowagi: $F=-kx$. Dla masy $m$ równanie $m\ddot x+kx=0$ ma rozwiązanie $x=A\cos(\omega_0t+\phi)$, gdzie $\omega_0=\sqrt{k/m}$, $A$ jest amplitudą, a $\phi$ fazą początkową. Okres wynosi $T=2\pi/\omega_0$. Energia $E=\tfrac12m\dot x^2+\tfrac12kx^2$ jest stała bez strat. Małe drgania dowolnego układu wokół stabilnej równowagi często są w przybliżeniu harmoniczne, gdy potencjał można rozwinąć do drugiego rzędu w wychyleniu.

Lepkie tłumienie modeluje się siłą $-b\dot x$, co daje $m\ddot x+b\dot x+kx=0$. Gdy $b<2\sqrt{mk}$, ruch jest niedotłumiony: amplituda maleje wykładniczo, a częstość drgań jest mniejsza niż $\omega_0$. Dla większego tłumienia układ wraca do równowagi bez oscylacji. Wymuszenie okresowe $F_0\cos\omega t$ prowadzi po zaniku stanu przejściowego do drgań ustalonych o częstości wymuszenia. Amplituda zależy od $\omega$ i tłumienia; dla słabego tłumienia osiąga maksimum w pobliżu częstości własnej, co nazywamy rezonansem. Zwiększenie tłumienia obniża i poszerza maksimum. Rezonans nie oznacza nieograniczonego wzrostu rzeczywistej amplitudy, ponieważ straty i nieliniowość ograniczają odpowiedź. Przykładami są zawieszenie samochodu oraz wymuszone drgania mostu; należy rozróżniać częstość własną układu i częstość zewnętrznego pobudzania.

Wykres wychylenia w czasie pozwala odróżnić amplitudę, okres i fazę, natomiast wykres energii pokazuje zamianę energii sprężystej na kinetyczną. W położeniu równowagi prędkość jest największa, a wychylenie i energia potencjalna najmniejsze; w skrajnych położeniach sytuacja jest odwrotna. Przy wymuszeniu ważna jest również różnica faz między siłą a odpowiedzią: w pobliżu rezonansu zmienia się ona gwałtownie wraz z częstością. W rzeczywistym zawieszeniu tłumienie ma być kompromisem między szybkim uspokojeniem drgań a utratą komfortu i nadmiernym usztywnieniem układu.

\subsection{6. Prawo powszechnego ciążenia, prawa Keplera i energia grawitacyjna}
\paragraph{Zakres.} Prawo Newtona, pole i potencjał grawitacyjny, orbity oraz empiryczne prawa Keplera.

Dwa punktowe ciała o masach $m_1$ i $m_2$, odległe o $r$, przyciągają się siłą o wartości $F=Gm_1m_2/r^2$ wzdłuż prostej łączącej ich środki; $G$ jest stałą grawitacji. Dla rozkładów o symetrii sferycznej na zewnątrz można traktować masę jako skupioną w środku. Pole grawitacyjne masy centralnej $M$ ma postać $\mathbf g=-GM\mathbf r/r^3$, czyli jest przyspieszeniem jednostkowej masy próbnej. Potencjał na jednostkę masy przy wyborze zera w nieskończoności wynosi $\varphi=-GM/r$, a energia potencjalna ciała $m$ to $U=m\varphi=-GMm/r$. Blisko powierzchni Ziemi, dla małych zmian wysokości, różnica energii przybliża się przez $mg\Delta h$.

Kepler opisał orbity planet: są elipsami z Słońcem w jednym ognisku; promień wodzący zakreśla równe pola w równych czasach; kwadrat okresu obiegu jest proporcjonalny do sześcianu wielkiej półosi. Dla masy planety pomijalnej wobec $M$ zachodzi $T^2=4\pi^2a^3/(GM)$. Drugie prawo Keplera jest skutkiem zachowania momentu pędu przy centralnej sile grawitacji. Dla orbity kołowej równowaga siły grawitacji i dośrodkowej daje $v=\sqrt{GM/r}$; całkowita energia mechaniczna wynosi wtedy $E=-GMm/(2r)$. Związek ten nie dotyczy dowolnej orbity ani punktu na elipsie jako energii kinetycznej. Model Newtonowski jest bardzo dobrym przybliżeniem dla słabych pól i prędkości znacznie mniejszych od prędkości światła; w silnym polu potrzebna jest ogólna teoria względności.

Potencjał grawitacyjny jest skalarem, co często upraszcza obliczenia w porównaniu z sumowaniem wektorowych sił. Dla kilku źródeł potencjały dodają się, ponieważ teoria Newtona jest liniowa w tym zakresie. Ujemny znak energii związanej orbity oznacza, że trzeba dostarczyć energię, aby oddalić ciało do nieskończoności z prędkością końcową równą zeru. W drugiej zasadzie Keplera można też odczytać zmienność prędkości na elipsie: planeta porusza się szybciej w pobliżu Słońca, a wolniej daleko od niego. Związek okresu i wielkiej półosi pozwala wyznaczać masy ciał centralnych z obserwacji orbit.

\subsection{7. Hydrostatyka i hydrodynamika płynów}
\paragraph{Zakres.} Ciśnienie, prawo Pascala i Archimedesa, równanie ciągłości oraz równanie Bernoulliego.

Ciśnienie jest siłą prostopadłą przypadającą na jednostkę powierzchni: $p=F_\perp/A$; w płynie w spoczynku działa jednakowo we wszystkich kierunkach w danym punkcie. W jednorodnym płynie o gęstości $\rho$ ciśnienie rośnie z głębokością: $p=p_0+\rho gh$, gdzie $p_0$ jest ciśnieniem na powierzchni, a $h$ głębokością. Prawo Pascala stwierdza, że przyrost ciśnienia w zamkniętym płynie przenosi się bez osłabienia; w prasie hydraulicznej $F_1/A_1=F_2/A_2$, kosztem większego przesunięcia mniejszego tłoka. Zasada Archimedesa mówi, że siła wyporu ma wartość ciężaru wypartego płynu, $F_w=\rho_{\rm pł}gV_{\rm zan}$. Ciało pływa, gdy równoważy ciężar, a nie dlatego, że jego ciężar znika.

Dla ustalonego przepływu nieściśliwego płynu przez przewód bilans masy daje równanie ciągłości $A_1v_1=A_2v_2$, gdzie $A$ jest polem przekroju, a $v$ średnią prędkością. Dla płynu idealnego, nielepkiego, o stałej gęstości, wzdłuż jednej linii prądu obowiązuje Bernoulli:
\[
p+\frac12\rho v^2+\rho gz=\text{const}.
\]
Składniki oznaczają odpowiednio energię ciśnienia, kinetyczną i potencjalną na jednostkę objętości. Zwężenie poziomej rury zwiększa prędkość i zmniejsza ciśnienie statyczne w modelu idealnym. Rzeczywiste płyny mają lepkość; straty energii, turbulencja, pompy i różnice między liniami prądu wymagają uzupełnienia równania. Bernoulliego nie należy stosować bez sprawdzenia założeń, na przykład do gwałtownie zmiennego przepływu z dużymi stratami.

Ciśnienie hydrostatyczne zależy od głębokości, a nie od kształtu naczynia; dlatego naczynia połączone z tym samym płynem osiągają zgodne poziomy przy jednakowym ciśnieniu nad powierzchnią. Wypór wynika z różnicy ciśnień działających na dolną i górną część zanurzonego ciała. W przepływie należy odróżniać ciśnienie statyczne od całkowitego bilansu energii płynu. Przed użyciem równania ciągłości warto ustalić, czy ciecz można uznać za nieściśliwą i czy przepływ jest ustalony. Przy znacznej lepkości część energii mechanicznej zamienia się w energię wewnętrzną, więc idealny bilans Bernoulliego nie jest samodzielnie wystarczający.

\subsection{8. Transformacje Galileusza i Lorentza, kinematyka i dynamika relatywistyczna}
\paragraph{Zakres.} Klasyczne i relatywistyczne transformacje układów, dylatacja czasu, skrócenie długości, pęd i energia kinetyczna.

Transformacja Galileusza zakłada absolutny czas: dla układów względnie poruszających się z prędkością $V$ wzdłuż osi $x$, $x'=x-Vt$ oraz $t'=t$. Jest poprawna, gdy $V\ll c$. Szczególna teoria względności przyjmuje jednakowe prawa fizyki w układach inercjalnych i niezmienniczość prędkości światła $c$. Wzdłuż osi ruchu transformacja Lorentza jest
\[
x'=\gamma(x-Vt),\qquad t'=\gamma(t-Vx/c^2),\qquad
\gamma=\frac{1}{\sqrt{1-V^2/c^2}}.
\]
Własny czas $\Delta\tau$ mierzony zegarem poruszającym się wraz z obiektem jest krótszy niż odstęp czasu $\Delta t=\gamma\Delta\tau$ w układzie, względem którego obiekt się porusza. Długość $L_0$ mierzona w układzie spoczynkowym pręta jest długością własną; w kierunku ruchu obserwuje się $L=L_0/\gamma$. Porównania wymagają równoczesnych pomiarów końców pręta w danym układzie.

Relatywistyczny pęd wynosi $\mathbf p=\gamma m\mathbf v$, gdzie $m$ jest masą spoczynkową. Energia całkowita $E=\gamma mc^2$, energia spoczynkowa $mc^2$, a energia kinetyczna $K=(\gamma-1)mc^2$. Dla $v\ll c$ rozwinięcie daje $K\approx\tfrac12mv^2$, więc odzyskujemy mechanikę klasyczną. Względność jednoczesności wyjaśnia, dlaczego transformacja czasu zawiera $x$; nie jest to jedynie poprawka zegara. Zasada zachowania energii-pędu ma relatywistycznie spójną postać, a niezmiennik spełnia $E^2=p^2c^2+m^2c^4$. Dla obiektów bezmasowych $E=pc$. Wzory klasyczne nie powinny być używane przy prędkościach porównywalnych z $c$.

Nie można traktować dylatacji czasu jako błędu zegara: obserwatorzy w różnych układach inercjalnych przypisują temu samemu zestawowi zdarzeń różne odstępy czasu, zgodnie z transformacją Lorentza. Każdy obserwator mierzy jednak tę samą prędkość światła w próżni. W granicy małych prędkości współczynnik gamma zbliża się do jedności, a składnik mieszający położenie z czasem staje się pomijalny, co uzasadnia przejście do opisu Galileuszowego. Dla porównania wielkości należy zawsze określić układ pomiaru; nie istnieje jeden uprzywilejowany układ spoczynkowy dla ruchu jednostajnego.

\subsection{9. Gaz doskonały i przemiany gazowe}
\paragraph{Zakres.} Równanie stanu gazu doskonałego, interpretacja mikroskopowa i podstawowe przemiany termodynamiczne.

Gaz doskonały jest modelem rozrzedzonego gazu, w którym cząsteczki traktuje się jako punktowe, ich wzajemne oddziaływania pomija poza krótkimi zderzeniami, a zderzenia są sprężyste. Równanie stanu brzmi $pV=nRT=Nk_{\rm B}T$, gdzie $p$ jest ciśnieniem, $V$ objętością, $n$ liczbą moli, $N$ liczbą cząsteczek, $R$ stałą gazową, $k_{\rm B}$ stałą Boltzmanna, a $T$ temperaturą bezwzględną. Dla ustalonej ilości gazu $pV/T$ jest stałe. Temperatura wiąże się ze średnią energią translacyjną cząsteczki: $\langle K_{\rm tr}\rangle=\tfrac32k_{\rm B}T$.

Przemiana izotermiczna zachodzi przy stałym $T$, zatem dla gazu doskonałego $pV=\text{const}$. Przemiana izochoryczna ma stałe $V$ i pracę objętościową równą zero; ciśnienie jest proporcjonalne do temperatury. W izobarycznej $p$ jest stałe, więc praca gazu przy zmianie objętości wynosi $W=p\Delta V$. Przemiana adiabatyczna nie wymienia ciepła z otoczeniem, $Q=0$; dla odwracalnej przemiany gazu doskonałego o stałych pojemnościach cieplnych $pV^\gamma=\text{const}$, gdzie $\gamma=C_p/C_V$. Adiabatyczna nie znaczy automatycznie izentropowa: izolowany układ z tarciem może produkować entropię. Wykres $p(V)$ pozwala odczytać pracę jako pole pod krzywą. Model gazu doskonałego zawodzi przy dużej gęstości, niskiej temperaturze i blisko skraplania, gdy objętość cząsteczek i oddziaływania przestają być pomijalne.

W przemianie trzeba odróżnić warunek definiujący jej przebieg od wielkości, które mogą się zmieniać. Na przykład stała temperatura nie oznacza braku przekazywania energii: podczas izotermicznego rozprężania gaz może pobierać ciepło, które następnie zamienia na pracę. Na wykresie ciśnienia od objętości izoterma gazu doskonałego ma kształt hiperboli, a adiabata opada szybciej, jeśli pojemności cieplne są dodatnie. Ilość gazu należy wyrażać spójnie jako liczbę moli albo liczbę cząsteczek i dobrać odpowiednią stałą. Temperatura w równaniu stanu jest bezwzględna, więc użycie skali Celsjusza prowadziłoby do błędnego wyniku.

\subsection{10. Pierwsza zasada termodynamiki i zasada ekwipartycji}
\paragraph{Zakres.} Bilans energii wewnętrznej, praca i ciepło, pojemności cieplne oraz ekwipartycja energii.

Pierwsza zasada termodynamiki jest zasadą zachowania energii dla układu makroskopowego. Przyjmując, że $W$ oznacza pracę wykonaną przez układ, zapisujemy $\Delta U=Q-W$, gdzie $U$ jest energią wewnętrzną, a $Q$ ciepłem dostarczonym układowi. Ciepło i praca są sposobami przekazywania energii, a nie wielkościami zawartymi w stanie układu. Dla quasi-statycznej pracy objętościowej $W=\int p\,dV$. W cyklu, po powrocie do stanu początkowego, $\Delta U=0$, dlatego całe dostarczone ciepło jest równe pracy wykonanej przez układ.

Dla gazu doskonałego wewnętrzna energia zależy tylko od temperatury. Dla $n$ moli $dU=nC_VdT$, gdzie $C_V$ jest molową pojemnością cieplną przy stałej objętości. W przemianie izochorycznej $W=0$, więc $Q=\Delta U$; w izotermicznej gazu doskonałego $\Delta U=0$, więc dostarczone ciepło równa się pracy. Zasada ekwipartycji mówi, że w klasycznej równowadze każdy niezależny kwadratowy składnik energii ma średnią wartość $\tfrac12 k_{\rm B}T$ na cząsteczkę. Jednoatomowy gaz idealny ma trzy translacyjne stopnie swobody, stąd $U=\tfrac32Nk_{\rm B}T$ i $C_V=\tfrac32R$ na mol. Wibracyjne stopnie swobody mogą być „zamrożone” kwantowo w niskiej temperaturze; ekwipartycja nie jest uniwersalna poza klasycznym reżimem. Warto zawsze jawnie podać konwencję znaków pracy, bo niektóre podręczniki definiują pracę wykonaną nad układem.

Przy analizie należy wskazać granicę układu i konwencję znaków, a następnie wypisać dopływy i odpływy ciepła oraz pracy. Dla układu zamkniętego energia wewnętrzna obejmuje mikroskopową energię ruchu i oddziaływań cząsteczek, a nie energię ruchu całego układu jako całości. Pojemność cieplna zależy od ograniczenia przemiany: przy stałym ciśnieniu dostarczone ciepło obejmuje również pracę rozprężania, dlatego molowa pojemność $C_p$ jest większa niż $C_V$ dla gazu doskonałego. Klasyczna ekwipartycja dobrze opisuje wiele gazów w temperaturach umiarkowanych, ale nie zastępuje eksperymentalnych wartości pojemności w każdej substancji.

\subsection{11. Druga zasada termodynamiki, ciepło zredukowane i entropia}
\paragraph{Zakres.} Kierunek przemian, sformułowania drugiej zasady, entropia i nierówność Clausiusa.

Pierwsza zasada bilansuje energię, ale nie wyznacza kierunku procesów. Druga zasada stwierdza między innymi, że nie istnieje cykliczna maszyna, której jedynym skutkiem byłoby pobieranie ciepła z jednego zbiornika i całkowite zamienianie go na pracę. Ciepło samorzutnie przepływa od ciała cieplejszego do chłodniejszego, nie odwrotnie bez dodatkowych skutków. Dla procesu odwracalnego entropię definiuje się przez $dS=\delta Q_{\rm rev}/T$, a ciepło zredukowane to $\delta Q_{\rm rev}/T$; $\delta Q$ jest różniczką niezupełną, bo zależy od drogi, podczas gdy $S$ jest funkcją stanu.

Dla dowolnego cyklu obowiązuje nierówność Clausiusa $\oint \delta Q/T\leq0$, z równością w granicy odwracalnej. Entropia układu izolowanego nie maleje: $\Delta S_{\rm izol}\geq0$. W procesie nieodwracalnym, takim jak swobodne rozprężanie gazu, entropia rośnie mimo braku wymiany ciepła z otoczeniem. Dla wymiany ciepła $Q$ ze zbiornikiem o stałej temperaturze $T$ zmiana entropii zbiornika wynosi $-Q/T$, jeśli $Q$ wpływa do badanego układu. Odwracalność jest idealizacją wymagającą braku tarcia, skończonych różnic temperatur i innych źródeł produkcji entropii. Sprawność silnika cieplnego między temperaturami $T_h$ i $T_c$ nie może przekraczać granicy Carnota $\eta_{\max}=1-T_c/T_h$ (temperatury w kelwinach). W praktyce nie należy utożsamiać entropii wyłącznie z „nieporządkiem”; jest precyzyjną wielkością termodynamiczną, która ilościowo ogranicza możliwe przemiany.

Entropię można obliczać wzdłuż dowolnej wygodnej drogi odwracalnej łączącej te same stany, nawet gdy rzeczywista przemiana przebiegła nieodwracalnie, ponieważ jej zmiana zależy tylko od stanów końcowych. Dla ciała o stałej pojemności cieplnej ogrzanego od temperatury $T_1$ do $T_2$ zmiana entropii jest dodatnia, jeśli temperatura końcowa jest większa. W bilansie całego procesu trzeba uwzględnić zarówno układ, jak i otoczenie; lokalny spadek entropii nie przeczy zasadzie, jeśli wzrost w otoczeniu jest większy. Ta perspektywa pozwala rozumieć, dlaczego silniki i lodówki wymagają wymiany energii z więcej niż jednym zbiornikiem oraz dlaczego rzeczywiste urządzenia mają ograniczoną sprawność.

\subsection{12. Kohezja, napięcie powierzchniowe, kapilarność i ruchy Browna}
\paragraph{Zakres.} Oddziaływania międzycząsteczkowe, energia powierzchniowa, zjawiska kapilarne oraz mikroskopowy ruch cząstek.

Kohezja oznacza przyciąganie cząsteczek tej samej substancji, a adhezja oddziaływanie między różnymi substancjami. Cząsteczka we wnętrzu cieczy jest otoczona sąsiadami ze wszystkich stron; cząsteczka przy powierzchni ma niezrównoważone oddziaływania. Zwiększenie powierzchni wymaga pracy, co opisuje napięcie powierzchniowe $\sigma=dW/dA$, mierzone w niutonach na metr, równoważnie dżulach na metr kwadratowy. Dla błonki z dwiema powierzchniami trzeba uwzględnić obie. Napięcie dąży do zmniejszenia powierzchni, dlatego krople są zbliżone do kulistych. Różnica ciśnień po obu stronach zakrzywionej powierzchni wynika z prawa Laplace’a; dla kulistej kropli cieczy $\Delta p=2\sigma/R$.

W cienkiej kapilarze równowaga składowej pionowej siły napięcia powierzchniowego z ciężarem słupa cieczy daje wysokość wzniesienia
\[
h=\frac{2\sigma\cos\theta}{\rho gr},
\]
gdzie $r$ jest promieniem rurki, $\theta$ kątem zwilżania, a $\rho$ gęstością cieczy. Dodatnie $h$ oznacza wznoszenie, ujemne obniżenie; wzór zakłada cienką rurkę, statyczną ciecz i pomijalny wpływ menisku na objętość. Ruchy Browna to nieregularne ruchy zawieszonych drobin wywołane zderzeniami z cząsteczkami płynu. W jednowymiarowym dyfuzji średni kwadrat przesunięcia wynosi $\langle x^2\rangle=2Dt$, gdzie $D$ jest współczynnikiem dyfuzji. Zjawisko stanowi dowód kinetyczno-molekularnej natury materii; nie wynika z celowego, uporządkowanego popychania drobiny przez płyn. Ruch jest bardziej widoczny dla małych cząstek, w cieplejszej cieczy i przy mniejszej lepkości.

W doświadczeniu z błonką cieczy zwiększenie jej powierzchni wymaga pracy, co obrazuje mechaniczne znaczenie napięcia powierzchniowego. Kąt zwilżania odróżnia ciecz rozpływającą się po powierzchni od tworzącej krople; sam znak wysokości kapilarnej nie wystarcza do określenia kształtu menisku bez znajomości zwilżania. Zależność od promienia rurki przewiduje wyraźnie większe wzniesienie w węższej kapilarze, o ile można pominąć wpływ zanieczyszczeń i nieregularności ścian. Ruchy Browna są statystyczne, więc opisuje się je średnimi i rozkładami, a nie pojedynczą przewidywalną trajektorią. Zjawiska powierzchniowe mają szczególne znaczenie dla małych kropel i porowatych materiałów.

\subsection{13. Zjawiska transportowe w gazach}
\paragraph{Zakres.} Dyfuzja, przewodnictwo cieplne, lepkość i przewodnictwo elektryczne w ujęciu makroskopowym i kinetycznym.

Zjawiska transportowe powstają, gdy w gazie istnieją nierównomierności stężenia, temperatury, prędkości makroskopowej lub potencjału elektrycznego. Cząsteczki przenoszą wtedy masę, energię, pęd albo ładunek, dążąc do wyrównania wielkości makroskopowych. Prawo Ficka dla strumienia cząstek ma postać $\mathbf j_N=-D\nabla n$, gdzie $n$ jest koncentracją, a $D$ współczynnikiem dyfuzji. Strumień ciepła spełnia prawo Fouriera $\mathbf q=-\kappa\nabla T$, gdzie $\kappa$ jest przewodnością cieplną. Ujemne znaki oznaczają transport w kierunku malejącego stężenia lub temperatury.

Lepkość opisuje transport pędu między warstwami płynu o różnych prędkościach. Dla przepływu newtonowskiego naprężenie ścinające jest proporcjonalne do gradientu prędkości; współczynnikiem proporcjonalności jest lepkość dynamiczna $\eta$. W gazie rozrzedzonym wszystkie te współczynniki można rozumieć kinetycznie przez przebywanie cząsteczek między zderzeniami, ich prędkości termiczne i średnią drogę swobodną. W przybliżeniu $D\sim \bar v\lambda/3$, gdzie $\bar v$ jest charakterystyczną prędkością, a $\lambda$ średnią drogą swobodną; analogicznie przewodnictwo i lepkość zależą od transportu energii i pędu. Przewodnictwo elektryczne opisuje prawo $\mathbf j=\sigma_{\rm el}\mathbf E$ dla liniowego, izotropowego ośrodka, gdzie $\mathbf j$ jest gęstością prądu, $\mathbf E$ polem, a $\sigma_{\rm el}$ przewodnością. W gazie nośnikami są naładowane cząstki; zwykły gaz obojętny słabo przewodzi, natomiast gaz zjonizowany może dobrze przewodzić. Prawa fenomenologiczne zakładają lokalną równowagę i małe gradienty; przy silnym rozrzedzeniu albo bardzo szybkich zmianach konieczny jest opis kinetyczny.

Wspólnym elementem tych zjawisk jest transport przez chaotyczny ruch cząsteczek, choć transportowana wielkość jest inna w każdym przypadku. Gradient temperatury wywołuje strumień energii, gradient stężenia strumień cząstek, a różnica prędkości warstw strumień pędu. Współczynniki transportu są dodatnie dla materiału pasywnego, co sprawia, że strumienie przeciwdziałają nierównomiernościom. W gazie przy większym ciśnieniu zderzenia są częstsze i droga swobodna krótsza, lecz zależności współczynników od ciśnienia zależą od reżimu i przybliżeń kinetycznych. Przewodnictwa elektrycznego nie należy utożsamiać z dyfuzją obojętnych cząsteczek, ponieważ wymaga mobilnych nośników ładunku i pola elektrycznego.

\subsection{14. Potencjał elektryczny, natężenie pola, strumień i prawo Gaussa}
\paragraph{Zakres.} Pole i potencjał elektrostatyczny, związek między nimi, strumień oraz prawo Gaussa.

Pole elektryczne definiuje się jako siłę działającą na dodatni ładunek próbny podzieloną przez ten ładunek: $\mathbf E=\mathbf F/q$. Pole punktowego ładunku $Q$ w próżni wynosi $\mathbf E=(1/4\pi\varepsilon_0)(Q/r^2)\hat{\mathbf r}$. W elektrostatyce pole jest zachowawcze, więc różnica potencjałów między punktami to ujemna praca pola na jednostkę ładunku: $V_B-V_A=-\int_A^B\mathbf E\cdot d\mathbf l$. Potencjał jest energią potencjalną na jednostkę ładunku, a $\mathbf E=-\nabla V$. Pole wskazuje kierunek najszybszego spadku potencjału; jego jednostką jest wolt na metr, równoważny niutonowi na kulomb.

Strumień pola przez powierzchnię $S$ jest $\Phi_E=\int_S\mathbf E\cdot d\mathbf A$, gdzie wektor elementu powierzchni jest skierowany na zewnątrz dla powierzchni zamkniętej. Prawo Gaussa w próżni mówi
\[
\oint_S\mathbf E\cdot d\mathbf A=\frac{Q_{\rm wew}}{\varepsilon_0}.
\]
Jest ono zawsze prawdziwe elektrostatycznie, ale szczególnie użyteczne obliczeniowo przy symetrii sferycznej, cylindrycznej lub płaskiej. Dla ładunku punktowego powierzchnia sferyczna daje prawo Coulomba. Ładunki na zewnątrz powierzchni mogą tworzyć pole na jej powierzchni, ale ich łączny strumień przez zamkniętą powierzchnię jest zerowy. W materii używa się także wektora indukcji elektrycznej $\mathbf D$ i ładunków swobodnych; w próżni $\mathbf D=\varepsilon_0\mathbf E$. Potencjał można wyznaczyć przez całkowanie pola lub z rozkładu ładunków, pamiętając, że jego poziom zerowy wybiera się umownie.

Przy kilku ładunkach pole sumuje się wektorowo, natomiast potencjał sumuje się algebraicznie, co czyni go często wygodniejszym w obliczeniach. Powierzchnie ekwipotencjalne są prostopadłe do linii pola elektrostatycznego, ponieważ wzdłuż takiej powierzchni praca pola jest zerowa. Prawo Gaussa nie mówi, że pole jest wszędzie równe zero, gdy całkowity ładunek wewnętrzny znika; oznacza jedynie zerowy strumień netto, a pole może przechodzić przez powierzchnię w obu kierunkach. Przy wyborze powierzchni Gaussa należy wykorzystać symetrię źródeł, a nie zakładać symetrię, której rozkład ładunku nie ma. W przewodniku w równowadze elektrostatycznej pole wewnątrz materiału jest zerowe.

\subsection{15. Polaryzacja dielektryków, elektrostrykcja, piezoelektryczność i ferroelektryczność}
\paragraph{Zakres.} Dielektryki w polu, polaryzacja, pojemność kondensatora oraz sprzężenia elektryczne z odkształceniem.

Dielektryk nie ma swobodnych nośników ładunku, które mogłyby przepływać przez cały materiał tak, jak w metalu, lecz pole elektryczne przesuwa ładunki związane i orientuje dipole. Polaryzacja $\mathbf P$ jest dipolowym momentem na jednostkę objętości. W liniowym, izotropowym dielektryku $\mathbf P=\varepsilon_0\chi_e\mathbf E$, gdzie $\chi_e$ jest podatnością elektryczną, a $\mathbf D=\varepsilon_0\mathbf E+\mathbf P=\varepsilon\mathbf E$. Wypełnienie kondensatora płaskiego dielektrykiem zwiększa pojemność z $C_0=\varepsilon_0A/d$ do $C=\varepsilon A/d=\varepsilon_rC_0$, jeśli materiał całkowicie wypełnia szczelinę i pole można uznać za jednorodne. Przy stałym ładunku zmniejsza to pole, a przy stałym napięciu źródło dostarcza dodatkowy ładunek; należy rozróżniać te warunki.

Elektrostrykcja to odkształcenie dielektryka związane z polaryzacją, na ogół zależne kwadratowo od pola i występujące także w materiałach centrosymetrycznych. Efekt piezoelektryczny oznacza sprzężenie: naprężenie może wywołać polaryzację, a pole może wywołać odkształcenie; wymaga odpowiedniej struktury krystalicznej bez środka symetrii. Ferroelektryki mają spontaniczną polaryzację, którą można przełączać polem; ich odpowiedź jest nieliniowa i wykazuje histerezę, więc proste równanie liniowe nie wystarcza. Przykładowo piezoelektryczny kryształ może przetwarzać drgania mechaniczne na sygnał elektryczny, a ferroelektryk służyć jako pamięć. Granice liniowego modelu wyznaczają silne pola, anizotropia, częstotliwość, temperatura i nasycenie polaryzacji. Energia kondensatora liniowego wynosi $U=\tfrac12CV^2=\tfrac12QV$.

Polaryzacja tworzy na powierzchni dielektryka ładunki związane, które modyfikują pole w materiale i na jego granicach. W jednorodnym dielektryku ich opis można sprowadzić do efektywnej przenikalności, lecz przy granicach warstw należy uwzględnić warunki brzegowe pola. Różne mechanizmy polaryzacji, takie jak przesunięcie elektronów i jonów albo obrót dipoli, mogą reagować z różną szybkością, dlatego odpowiedź zależy od częstotliwości pola. Elektrostrykcja jest zwykle parzysta względem odwrócenia pola, podczas gdy piezoelektryczność może zmieniać znak wraz z kierunkiem pola lub naprężenia. Rozróżnienie tych efektów pomaga prawidłowo interpretować odkształcenia materiału i wybierać elementy do czujników, przetworników oraz magazynowania ładunku.

\subsection{16. Klasyczne przewodnictwo elektryczne}
\paragraph{Zakres.} Prąd, gęstość prądu, model Drudego, opór oraz lokalna postać prawa Ohma.

Prąd elektryczny jest szybkością przepływu ładunku przez przekrój: $I=dQ/dt$. Umowny kierunek prądu jest kierunkiem ruchu ładunków dodatnich; w metalach elektrony przemieszczają się przeciwnie. Gęstość prądu $\mathbf j$ opisuje lokalny przepływ, a natężenie $I=\int_S\mathbf j\cdot d\mathbf A$. Równanie ciągłości ładunku ma postać $\partial\rho_q/\partial t+\nabla\cdot\mathbf j=0$, wyrażając zachowanie ładunku. W jednorodnym przewodniku o długości $L$ i przekroju $A$ opór definiuje $R=V/I$; dla materiału omowego $R=\rho_{\rm el}L/A$, gdzie $\rho_{\rm el}$ jest opornością właściwą.

W modelu Drudego elektrony przewodnictwa poruszają się między zderzeniami z jonami sieci. Pole nadaje im średnią prędkość dryfu, a zderzenia przekazują energię sieci i tworzą opór. W przybliżeniu $\mathbf j=\sigma\mathbf E$, gdzie $\sigma$ jest przewodnością; w prostym modelu $\sigma=ne^2\tau/m$, przy koncentracji elektronów $n$, ładunku elementarnym $e$, średnim czasie między zderzeniami $\tau$ i masie elektronu $m$. Prawo Ohma jest prawem materiałowym w określonym zakresie, a nie uniwersalnym prawem dla wszystkich elementów. Temperatura zwykle zwiększa drgania sieci metalicznej i opór metali, choć zależność różni się dla półprzewodników i innych materiałów. Moc cieplna wydzielana w oporze wynosi $P=IV=I^2R=V^2/R$ dla elementu omowego. Model klasyczny nie opisuje szczegółowo struktury pasmowej, nadprzewodnictwa ani wielu własności kwantowych metali, lecz przydatnie wprowadza mechanizm dryfu i rozpraszania.

Prędkość dryfu jest na ogół znacznie mniejsza od prędkości chaotycznego ruchu elektronów, mimo że zmiana pola w obwodzie może rozchodzić się szybko. Nie należy zatem wyobrażać sobie, że pojedynczy elektron przebywa cały obwód z szybkością sygnału elektrycznego. W stanie ustalonym prąd jest taki sam w kolejnych przekrojach przewodu szeregowego, choć lokalna gęstość prądu może zależeć od pola przekroju. Oporność właściwa jest cechą materiału w określonych warunkach, podczas gdy opór zależy również od geometrii próbki. Model Drudego wyjaśnia jakościowo te różnice, a bardziej dokładny opis wykorzystuje stany kwantowe elektronów w krysztale.

\subsection{17. Prawo Ohma, siła elektromotoryczna, prawa Kirchhoffa, praca i moc}
\paragraph{Zakres.} Obwody prądu stałego, źródła rzeczywiste, reguły Kirchhoffa oraz bilans energii i mocy.

Prawo Ohma dla opornika liniowego to $V=IR$, gdzie $V$ jest napięciem na elemencie, $I$ prądem, a $R$ jego oporem. Siła elektromotoryczna $\mathcal E$ nie jest siłą mechaniczną, lecz pracą źródła przypadającą na jednostkę ładunku, która napędza ładunki wbrew elektrostatycznemu spadkowi potencjału. Źródło o oporze wewnętrznym $r$ dostarczające prąd $I$ ma napięcie zaciskowe $V=\mathcal E-Ir$. Zatem przy zwarciu prąd jest ograniczony przez $r$ w modelu idealnym; rzeczywiste źródła mają dodatkowe ograniczenia i mogą się przegrzać.

Pierwsze prawo Kirchhoffa wynika z zachowania ładunku: algebraiczna suma prądów w węźle jest zerowa. Drugie wynika z bilansu energii i zachowawczości pola elektrostatycznego w obwodzie skupionym: suma zmian potencjałów i sił elektromotorycznych po zamkniętej pętli wynosi zero. Przy układaniu równań ustala się umowne kierunki prądów i obchodu; wynik ujemny wskazuje kierunek przeciwny. Rezystory szeregowe sumują opory, a równoległe sumują odwrotności oporów. Praca przekazana ładunkowi przez napięcie wynosi $W=qV$, a moc chwilowa $P=IV$. Dla opornika cała energia elektryczna zamienia się w ciepło z szybkością $I^2R$. Przy wielu źródłach i opornikach równania Kirchhoffa pozwalają wyznaczyć prądy, ale należy sprawdzić znaki oraz spełnienie bilansu mocy: moc dostarczana przez źródła równa się mocy pochłanianej przez elementy w stanie ustalonym. Przy wysokich częstotliwościach efekty indukcyjne i pojemnościowe wymagają bogatszego modelu niż czysto rezystancyjny.

W rozwiązaniu wielooczkowego obwodu przydatne jest konsekwentne oznaczenie prądów gałęziowych, biegunów źródeł i kierunku obchodzenia każdej pętli. Spadek napięcia na oporniku w kierunku umownego prądu ma przeciwny znak niż wzrost potencjału przez źródło od bieguna ujemnego do dodatniego. Jeśli obliczony prąd jest ujemny, równania pozostają poprawne, a rzeczywisty kierunek jest przeciwny do przyjętego. Moc źródła może być dodatnia lub ujemna zależnie od tego, czy dostarcza energię, czy jest ładowane. Te kontrole znaków pozwalają wykryć błędy rachunkowe bez zmieniania przyjętych kierunków w trakcie rozwiązania.

\subsection{18. Siła Ampère’a, oddziaływanie przewodów, indukcja i prawo Faradaya}
\paragraph{Zakres.} Siły magnetyczne na prądy, pole przewodów, strumień magnetyczny i indukcja elektromagnetyczna.

Na ładunek $q$ poruszający się z prędkością $\mathbf v$ w polu magnetycznym $\mathbf B$ działa składowa magnetyczna siły Lorentza $\mathbf F=q\mathbf v\times\mathbf B$. Na odcinek przewodnika z prądem działa siła Ampère’a $d\mathbf F=I\,d\mathbf l\times\mathbf B$. Dwa długie, równoległe przewody odległe o $r$, niosące prądy $I_1$ i $I_2$, oddziałują siłą na jednostkę długości $F/L=\mu_0I_1I_2/(2\pi r)$ w próżni: prądy zgodne przyciągają się, przeciwne odpychają. Kierunek można wyznaczyć regułą prawej dłoni. Siła magnetyczna na pojedynczy ładunek nie wykonuje pracy, ponieważ jest prostopadła do jego prędkości.

Strumień magnetyczny przez powierzchnię to $\Phi_B=\int_S\mathbf B\cdot d\mathbf A$. Prawo Faradaya-Lenza dla obwodu brzmi $\mathcal E_{\rm ind}=-d\Phi_B/dt$; dla cewki o $N$ zwojach, $\mathcal E=-N\,d\Phi_B/dt$. Znak minus wyraża regułę Lenza: prąd indukowany wytwarza pole przeciwdziałające zmianie strumienia, a niekoniecznie samemu polu. Strumień może się zmieniać wskutek zmiany pola, powierzchni lub kąta. W przypadku ruchomego przewodnika opis siły Lorentza i prawo Faradaya są zgodne. Przykładem jest generator, w którym obracająca się cewka zmienia strumień i indukuje napięcie. Indukcja własna cewki spełnia $\mathcal E_L=-L\,dI/dt$, a energia pola magnetycznego wynosi $\tfrac12LI^2$ dla liniowej indukcyjności. Prawo indukcji obowiązuje w pełnym opisie elektromagnetycznym; dla obwodów ruchomych trzeba starannie uwzględnić zarówno zmienność pola, jak i ruch przewodnika.

W zadaniach o indukcji warto najpierw ustalić orientację normalnej do powierzchni, a dopiero później wyznaczać znak strumienia i siły elektromotorycznej. Reguła Lenza jest konsekwencją zachowania energii: prąd nie może samorzutnie wzmacniać zmiany, która go wywołała, bez dopływu pracy z zewnątrz. Przesuwanie magnesu szybciej względem cewki zwiększa szybkość zmiany strumienia, a więc zwykle zwiększa wartość indukowanego napięcia. Z kolei przy stałym strumieniu nie ma indukcji, nawet jeśli samo pole jest silne. Przy wielu zwojach napięcia indukowane sumują się, jeśli ich orientacja jest zgodna.

\subsection{19. Materiały magnetyczne: diamagnetyki, paramagnetyki i ferromagnetyki}
\paragraph{Zakres.} Magnetyzacja, podatność magnetyczna i różnice fizyczne między trzema podstawowymi klasami materiałów.

Odpowiedź materii na pole magnetyczne opisuje magnetyzacja $\mathbf M$, czyli moment magnetyczny na jednostkę objętości. W prostym ośrodku liniowym i izotropowym $\mathbf M=\chi_m\mathbf H$, gdzie $\chi_m$ jest podatnością magnetyczną, a $\mathbf H$ natężeniem pola magnetycznego; indukcja spełnia $\mathbf B=\mu_0(\mathbf H+\mathbf M)$. W próżni $\mathbf M=0$ i $\mathbf B=\mu_0\mathbf H$. Zależności te są przybliżeniami: materiały anizotropowe, nieliniowe i histeretyczne wymagają bardziej ogólnego opisu.

Diamagnetyzm występuje zasadniczo we wszystkich materiałach: zmienne pole indukuje moment przeciwny do przyłożonego pola, więc podatność jest mała i ujemna, $\chi_m<0$. Paramagnetyki mają trwałe momenty atomowe, które pole częściowo ustawia; podatność jest mała i dodatnia, a uporządkowanie osłabia ruch cieplny. Ferromagnetyki, takie jak żelazo w odpowiednich warunkach, wykazują silne oddziaływanie między momentami, tworzenie domen i spontaniczne uporządkowanie poniżej temperatury Curie. Ich magnetyzacja zależy od historii pola: pętla histerezy wiąże się z namagnesowaniem szczątkowym i polem koercji. Powyżej temperatury Curie uporządkowanie ferromagnetyczne zanika, choć materiał może stać się paramagnetyczny. Miękkie materiały ferromagnetyczne łatwo magnesować i rozmagnesować, twarde zachowują silne namagnesowanie. Przy opisie trzeba odróżniać $\mathbf B$ i $\mathbf H$ oraz nie utożsamiać względnej przenikalności z uniwersalną stałą dla wszystkich pól i temperatur. Zastosowania obejmują rdzenie transformatorów, magnesy trwałe, czujniki i ekranowanie.

Klasyfikacja materiału zależy od jego odpowiedzi w zadanych warunkach, a nie wyłącznie od nazwy pierwiastka; temperatura, skład stopu i wcześniejsze magnesowanie mogą zmieniać obserwowane zachowanie. W diamagnetyku pole indukowane jest słabe, dlatego efekt często maskują silniejsze własności paramagnetyczne lub ferromagnetyczne. W ferromagnetyku domeny mogą być namagnesowane w różnych kierunkach, a przyłożone pole powoduje ich wzrost lub obrót, co prowadzi do nieliniowej krzywej magnesowania. Pole wewnątrz próbki zależy zarówno od odpowiedzi materiału, jak i kształtu próbki, który wpływa na pole rozmagnesowujące. W zastosowaniach inżynierskich znaczenie mają również straty histerezowe i prądy wirowe, które zamieniają część energii pola na ciepło.
\subsection{20. Drgania swobodne i wymuszone w obwodach LC i RLC; rezonans}
\paragraph{Zakres.} Rozpatruje się równania ruchu ładunku i napięcia w obwodach elektrycznych, zjawisko tłumienia oraz rezonans. Kluczowe są analogie do oscylatora mechanicznego oraz warunek maksymalnej odpowiedzi przy wymuszeniu harmonicznym. W obwodzie LC bez strat energia jest magazynowana w kondensatorze i w cewce: $E=\frac{q^2}{2C}+\frac{Li^2}{2}$, a dla prądu $i=\dot q$ równanie ma postać $L\ddot q+\frac{1}{C}q=0$. Stąd $\omega_0=1/\sqrt{LC}$ i rozwiązanie $q=q_0\cos(\omega_0 t+\phi)$. W obwodzie RLC przy zewnętrznym napięciu $u=U_0\cos\omega t$ równanie ruchu jest $L\ddot q+R\dot q+\frac{1}{C}q = U_0\cos\omega t$. Rozwiązanie składa się z części przejściowej i ustalonej; przy małym tłumieniu część przejściowa zanika, pozostaje drganie wymuszone o częstości $\omega$.\par
Dla obwodu szeregowym amplituda prądu jest $I_0=\frac{U_0}{\sqrt{R^2+(\omega L-1/\omega C)^2}}$. Maksimum występuje, gdy $\omega L=1/(\omega C)$, czyli $\omega=\omega_0$. Współczynnik dobroci $Q=\omega_0 L/R$ charakteryzuje ostrość rezonansu; przy $Q\gg1$ szerokość pasma wynosi $\Delta\omega\approx\omega_0/Q$. Przy rezonansie prąd i napięcie są w fazie dla rezystora, a odpowiedź prądu rośnie do maksimum. Zjawisko to jest wykorzystywane w filtrach, wzmacniaczach i obwodach selektywnych; w układach $LC$ i $RLC$ rezonans odpowiada dopasowaniu częstotliwości do częstotliwości własnej układu.\par
W praktyce istotna jest też energia: w obwodzie bez strat energia jest cyklicznie zamieniana między kondensatorem a cewką, a przy wymuszeniu źródło dostarcza energię równoważącą straty w rezystorze. Analogicznie do oscylatora mechanicznego, rezonans występuje przy zgodności częstości wymuszenia z częstością własną; wówczas amplituda i moc przenoszone do układu są największe.\par

Dla poprawnego opisu trzeba rozróżnić częstość własną nieobciążonego układu od częstości, przy której amplituda odpowiedzi osiąga maksimum. W układzie o skończonym oporze oscylacje swobodne zanikają wykładniczo, a ich częstość jest nieco mniejsza od wartości bezstratnej, jeśli tłumienie nie jest zaniedbywalne. Energia tracona w rezystorze w jednym okresie jest uzupełniana przez źródło podczas drgań ustalonych. Przy rezonansie napięcia na cewce i kondensatorze mogą być znacznie większe od napięcia generatora, lecz są przeciwne w fazie i ich wypadkowe napięcie reaktywne się znosi. Warto przy odpowiedzi ustnej podkreślić, że rezonans nie oznacza braku strat, tylko maksymalny transfer mocy lub odpowiedzi w danych warunkach. Model skupionych elementów zakłada, że rozmiary układu są małe w porównaniu z długością elektromagnetycznej fali.
\par

\subsection{21. Superpozycja fal, fale stojące i bieżące, dudnienia}
\paragraph{Zakres.} Rozważa się superpozycję fal, fale stojące i bieżące oraz dudnienia powstające przy nakładaniu się drgań o zbliżonych częstotliwościach. Kluczowe są faza, amplituda, długość fali i przepływ energii. Fala bieżąca w strunie ma postać $y_1=A\cos(kx-\omega t)$, a w przeciwnym kierunku $y_2=A\cos(kx+\omega t)$. Sumując, otrzymujemy $y=2A\cos kx\cos\omega t$, czyli falę stojącą: węzły występują dla $kx=(2n+1)\pi/2$, a antywęzły dla $kx=n\pi$. Odległość między sąsiednimi węzłami albo antywęzłami wynosi $\lambda/2$. Fala stojąca nie przenosi energii netto w kierunku propagacji, choć energia oscyluje między punktami o maksymalnej i minimalnej amplitudzie.\par
Dla dwóch fal o zbliżonych częstotliwościach $\omega_1$ i $\omega_2$ i tej samej amplitudzie $A$ wynikowa amplituda jest $2A\cos\left(\frac{\Delta\omega}{2}t\right)\cos\left(\frac{\omega_1+\omega_2}{2}t\right)$. To oznacza modulację amplitudy o częstości $|f_1-f_2|$, a w akustyce daje efekt dudnień: głośność zmienia się z częstotliwością beat. W falach mechanicznych w strunach i gazach prędkość zależy od napięcia i gęstości lub temperatury, a w falach elektromagnetycznych w próżni $v=c$. W instrumentach muzycznych dudnienia pozwalają na kalibrację tonów i sprawdzanie zgodności częstotliwości.\par
W jednowymiarowych układach zamkniętych lub otwartych pojawiają się stożkowe wzorce harmoniczne. Dla struny zamkniętej na końcach węzły są na końcach, a częstotliwości harmoniczne są $f_n = n v/(2L)$; dla rury otwartej na jednym końcu wzór przyjmuje $f_n=(2n+1)v/(4L)$. To pokazuje, że superpozycja rozwiązań normalnych prowadzi do wzorców stosowanych w akustyce i optyce.\par

Dudnienia należy odróżniać od fali stojącej: dudnienia wynikają z różnicy częstości, natomiast fala stojąca powstaje przy nakładaniu się fal o tej samej częstości biegnących w przeciwnych kierunkach. Węzły mają stale zerowe wychylenie, a w dudnieniach obwiednia zmienia się w czasie w każdym punkcie przestrzeni. Dla fal o nierównych amplitudach minimum natężenia nie spada do zera, a głębokość modulacji zależy od stosunku amplitud. Zasada superpozycji jest spełniona dla liniowych ośrodków i oznacza, że chwilowe wychylenia dodają się algebraicznie. W strunie szybkość fali zależy od napięcia i masy jednostki długości; zmiana tych parametrów zmienia długość fali przy zadanej częstości. Fale stojące są też podstawą rezonansów akustycznych, a ich wzory pozwalają wyznaczać parametry instrumentów.
\par

\subsection{22. Interferencja dwóch identycznych źródeł punktowych; maksima/minima i warunek pola dalekiego}
\paragraph{Zakres.} Rozpatruje się dwa identyczne źródła punktowe o tej samej częstotliwości i fazie, które emitują fale koherentne. Istotne są różnica dróg i warunek wystąpienia maksimów oraz minimów; omawia się również przybliżenie pola dalekiego (Fraunhofera). Przy punktach obserwacyjnych $P$ odległości do źródeł różnią się o $\Delta r = r_2-r_1$. Przy jednakowych amplitudach wynikowa amplituda jest $A_P=2A\cos(\Delta\phi/2)$, gdzie $\Delta\phi=k\Delta r$. Intensywność ma postać $I=4I_0\cos^2(\Delta\phi/2)$. Maksima pojawiają się dla $\Delta r = m\lambda$, minima dla $\Delta r=(m+1/2)\lambda$. W układzie dwóch źródeł odległych o $d$ i obserwowanych w kącie $\theta$ względem osi, przy małej odległości do ekranu $\Delta r\approx d\sin\theta$. Stąd $d\sin\theta=m\lambda$ dla maksimów i $d\sin\theta=(m+1/2)\lambda$ dla minimów.\par
Warunek pola dalekiego (Fraunhofera) ma postać $R\gg d^2/\lambda$, gdzie $R$ jest odległością źródeł do ekranu. Wtedy promienie z obu źródeł są praktycznie równoległe i można zastosować przybliżenie geometryczne. Dla małych kątów $\theta$ położenie prążka wynosi $y\approx R\theta$, a maksima pojawiają się przy $y_m\approx m\lambda R/d$. Rozstęp prążków jest $\Delta y\approx\lambda R/d$. Im większa odległość źródeł, tym większa różnica dróg i gęstszy układ prążków; im większa długość fali, tym większy odstęp. W praktyce ten model opisuje działanie wiązek laserowych, anten i eksperymentów interferencyjnych.\par

Różnica dróg jest dobrym sposobem opisu tylko wtedy, gdy fale są koherentne i zachowują ustaloną różnicę faz w czasie obserwacji. Jeżeli źródła emitują światło niespójne, prążki uśredniają się i obserwuje się sumę natężeń, a nie stabilny obraz interferencyjny. Przy różnych amplitudach maksima i minima pozostają w podobnych położeniach, lecz minima nie są całkowicie ciemne; dla amplitud $A_1$ i $A_2$ natężenie zależy od wyrazu interferencyjnego proporcjonalnego do $2A_1A_2\cos\Delta\phi$. W geometrii ekranowej wielkość prążków można zmierzyć i z niej wyznaczyć długość fali albo separację źródeł. Należy pamiętać, że przybliżenie małych kątów wymaga, aby obserwowany obszar ekranu był blisko osi optycznej. W doświadczeniu punktowe źródła można zastąpić dwiema szczelinami o wspólnym źródle oświetlenia.
\par

\subsection{23. Doświadczenie Younga, natężenie prążków; interferometr Michelsona}
\paragraph{Zakres.} Rozważamy układ dwóch szczelin w doświadczeniu Younga oraz zastosowanie interferometru Michelsona do pomiarów drobnych różnic dróg optycznych. W doświadczeniu Younga oświetlenie jest monochromatyczne, szczeliny są równo odległe, a ekran znajduje się w odległości $L$. Dla szczelin odległych o $d$ różnica faz wynosi $\Delta\phi = k d\sin\theta$, więc natężenie ma rozkład $I(\theta)=4I_1\cos^2\left(\frac{kd\sin\theta}{2}\right)$. Dla małych kątów, gdy $\sin\theta\approx y/L$, otrzymujemy $I(y)=4I_1\cos^2\left(\frac{\pi d}{\lambda L}y\right)$. Prążki jasne pojawiają się dla $d\sin\theta=m\lambda$, ciemne dla $d\sin\theta=(m+1/2)\lambda$, a odległość między sąsiednimi maksymalnymi wynosi $\Delta y = \lambda L/d$.\par
W praktyce na rozkład interferencyjny nakłada się dyfrakcja z każdej szczeliny; jeśli szczeliny mają skończoną szerokość, wzorzec obejmuje obwiednię dyfrakcyjną. To pokazuje, że światło nie jest prostą „promienistą” wiązką, ale falą, która przechodzi przez obie szczeliny i interferuje. Interferometr Michelsona ma półprzepuszczającą płytkę i dwa zwierciadła; jedna wiązka biegnie do zwierciadła ruchomego, druga do stałego, a po powrocie interferują. Zmiana długości drogi $\Delta$ daje przesunięcie fazy $\delta = 2\pi\Delta/\lambda$, a przy przesunięciu zwierciadła o $\lambda/2$ następuje przejście z maksimum do minimum. Stąd instrument mierzy bardzo małe odległości i zmiany indeksu załamania.\par

Doświadczenie Younga ma znaczenie historyczne, ponieważ wykazało interferencyjny charakter światła i umożliwiło precyzyjny pomiar długości fali. Warunkiem widocznych prążków jest spójność przestrzenna oświetlenia szczelin oraz dostatecznie wąskie pasmo widmowe źródła. Gdy źródło ma wiele długości fal, każda tworzy własny układ prążków, który nakłada się na pozostałe i ogranicza kontrast dalekich rzędów. W interferometrze Michelsona ruch lustra zmienia drogę optyczną dwukrotnie, ponieważ wiązka przebywa drogę do zwierciadła i z powrotem. Dlatego liczba przejść prążków przez punkt odniesienia pozwala obliczyć przesunięcie luster z dużą czułością. Wstawienie próbki do jednej gałęzi zmienia drogę optyczną przez zmianę współczynnika załamania, co umożliwia badanie gazów i materiałów. Stabilność mechaniczna układu jest niezbędna, ponieważ drgania mogą przesuwać prążki.
\par

\subsection{24. Efekt Dopplera dla fal mechanicznych i elektromagnetycznych oraz fala uderzeniowa}
\paragraph{Zakres.} Omawia się zmianę częstotliwości fal ze względu na ruch źródła i obserwatora oraz zjawisko fali uderzeniowej. Dla fal mechanicznych w ośrodku nieruchomym częstotliwość obserwowana wynosi $f' = f\frac{v\pm v_O}{v\mp v_S}$, gdzie $v$ jest prędkością fali w ośrodku, $v_O$ prędkością obserwatora, a $v_S$ prędkością źródła. Przy zbliżaniu się obserwatora do źródła częstotliwość rośnie; przy oddalaniu — maleje. Dla źródła zbliżającego się do obserwatora odstępy między frontami fal maleją, co odpowiada wzrostowi częstotliwości. W akustyce przy $v\approx 343\,\mathrm{m/s}$ są to efekty bardzo wyraźne, w astronomii zaś ocenia się ruch gwiazd i galaktyk.\par
W falach elektromagnetycznych w próżni związek między częstotliwością a długością fali jest $c=\lambda f$, a efekt Dopplera ma postać $\frac{\Delta\lambda}{\lambda}\approx \frac{v}{c}$ dla małych prędkości. W mechanice relatywistycznej trzeba użyć transformacji Lorentza; dla źródła lub obserwatora poruszającego się z prędkością $v$ względem wybranego układu należy uwzględnić czynnik $\gamma$ i kąt obserwacji. W astronomii przesunięcie ku czerwieni lub ku niebieskim indeksuje ruch obiektów kosmicznych.\par
Fala uderzeniowa powstaje, gdy źródło porusza się szybciej niż fala w ośrodku: $v_S>v$. Wtedy fronty fal są „nagradzane” i tworzą stożek Macha, dla którego $\sin\theta=v/v_S$. W akustyce prowadzi to do „przebicia dźwiękowego” i gwałtownego skoku ciśnienia; w aerodynamice i czasie wybuchów zjawisko to jest kluczowe dla wyznaczenia kształtu stożka uderzeniowego i siły działania.\par

W akustycznym efekcie Dopplera ośrodek wyznacza uprzywilejowany układ odniesienia: ruch źródła zmienia odstęp między kolejnymi frontami, a ruch obserwatora zmienia częstość ich napotykania. Z tego powodu prędkości źródła i obserwatora nie można w prosty sposób zastąpić jedną prędkością względną, gdy ośrodek jest istotny. Dla światła w próżni nie ma ośrodka spoczynkowego, a wynik zależy od względnego ruchu źródła i obserwatora oraz kierunku obserwacji. Przy dużych prędkościach klasyczne przybliżenie przesunięcia długości fali przestaje być wystarczające i należy stosować dokładny wzór relatywistyczny. Stożek Macha opisuje granicę zaburzeń, które nakładają się, gdy źródło przekracza prędkość fal; stosunek prędkości określa jego kąt. Przykładem są grzmot za szybkim samolotem oraz fale uderzeniowe w przepływie gazu.
\par

\subsection{25. Widmo elektromagnetyczne}
\paragraph{Zakres.} Opisuje się zakresy częstotliwości i długości fal promieniowania elektromagnetycznego, ich źródła, właściwości i zastosowania. Zależność $c=\lambda f$ i energia fotonu $E=h\nu$ są podstawowe. Widmo obejmuje fale radiowe, mikrofale, podczerwień, światło widzialne, nadfiolet, promieniowanie rentgenowskie i gamma. Długości fal w świetle widzialnym mieszczą się w przybliżeniu od $400\,\mathrm{nm}$ (fiolet) do $700\,\mathrm{nm}$ (czerwony); energia fotonu jest proporcjonalna do częstotliwości, zatem promieniowanie o krótszej długości fali ma większą energię.\par
Zakresy fal różnią się z powodu odmiennych mechanizmów emisji: radio i mikrofale związane z oscylacjami prądów oraz rotacją cząsteczek, podczerwień — z drganiami cząsteczkowymi, światło widzialne i UV z przejść elektronowych w atomach, rentgenowskie i gamma z przejść w jądrze lub z elektronów z głębokich powłok. Podział nie jest czysto arbitralny: istnieją obszary nakładające się i przejścia między nimi, związane z różnymi metodami generacji i detekcji. W praktyce widmo służy do komunikacji, obrazowania, spektroskopii, medycyny i kontroli jakości.\par
Promieniowanie jonizujące, zwłaszcza UV, X i gamma, ma energię wystarczającą do uszkadzania biomolekuł przez jonizację. Promieniowanie niejonizujące, takie jak mikrofale i podczerwień, zwykle nie rozrywa wiązań chemicznych, ale może prowadzić do ogrzewania lub wzbudzeń oscylacyjnych. Kluczowe jest rozróżnienie między energią fotonu, intensywnością wiązki i możliwością absorpcji w materiale; to właśnie w tych mechanizmach leży podstawa spektroskopii i diagnostyki.\par

Granice poszczególnych zakresów widma są umowne, ponieważ klasyfikacja opiera się przede wszystkim na długości fali lub energii fotonów, a nie na ostrych przejściach fizycznych. Promieniowanie wszystkich zakresów w próżni rozchodzi się z tą samą prędkością światła i podlega prawom Maxwella; różne sposoby detekcji wynikają z odmiennej energii i sprzężenia z materią. Fale radiowe wykrywa się antenami, światło detektorami fotoelektrycznymi, a promieniowanie jonizujące licznikami i detektorami półprzewodnikowymi. W spektroskopii położenie linii widmowych ujawnia strukturę energetyczną źródła, natomiast natężenie mówi o liczbie emitowanych lub pochłanianych fotonów. Dla fotonu światła widzialnego energia wynosi kilka elektronowoltów, a dla promieniowania gamma może być o wiele rzędów większa. Rozróżnienie to ma znaczenie w zastosowaniach, ochronie radiologicznej i badaniu kosmosu.
\par

\subsection{26. Huygens-Fresnel, dyfrakcja Fraunhofera w pojedynczej szczelinie i siatce dyfrakcyjnej}
\paragraph{Zakres.} Rozważa się zasadę Huygensa-Fresnela, dyfrakcję w pojedynczej szczelinie oraz działanie siatki dyfrakcyjnej; kluczowe są geometria układu i warunki tworzenia maksimów i minimów. Zasada Huygensa-Fresnela mówi, że każdy punkt frontu falowego jest źródłem wtórnych fal kulistych; ich suma amplitud w danym punkcie daje pole. Dla szczeliny szerokości $a$ w polu dalekim, intensywność ma postać $I(\theta)=I_0\left(\frac{\sin\beta}{\beta}\right)^2$, gdzie $\beta=(\pi a\sin\theta)/\lambda$. Minima pojawiają się przy $a\sin\theta = m\lambda$ dla $m=\pm1,\pm2,\ldots$; maksima główne leżą między nimi. W praktyce rozkład intensywności wytwarza szerokie maksimum centralne i symetryczne mniejsze maksima po bokach.\par
Siatka dyfrakcyjna składa się z wielu równoległych szczelin; jej rozkład intensywności jest iloczynem dyfrakcji pojedynczej szczeliny i interferencji między szczelinami. Warunek maksimów siatki: $d\sin\theta = m\lambda$, gdzie $d$ jest stałą siatki, a $m$ — rzędem widma. Rozdzielczość siatki jest $R=mN$ dla $N$ szczelin; większa liczba szczelin daje ostrzejsze maksima. Zjawisko to wykorzystuje się w spektrometrach do rozdzielania długości fal i analizy składu światła.\par
Przybliżenie Fraunhofera zakłada, że ekran jest bardzo daleko od przeszkody względem rozmiarów obiektu i długości fali. Wtedy promienie są prawie równoległe, a rozkład amplitudy można wyznaczyć jako transformację Fouriera wejścia. Dyfrakcja jest dowodem, że światło ma własności falowe, nie tylko promieniste. Jej skutki są odczuwalne w mikroskopii, optyce telekomunikacyjnej i badaniu struktur atomowych.\par

Zasada Huygensa-Fresnela nie oznacza, że wtórne fale należy traktować jako niezależne, bo ich amplitudy sumuje się z uwzględnieniem fazy i czynników kierunkowych. Dla pojedynczej szczeliny minima wynikają z kasowania wkładów z różnych obszarów apertury; przy pierwszym minimum można podzielić szczelinę na dwie części, których odpowiadające sobie punkty dają fale w przeciwfazie. Szerokość maksimum centralnego jest odwrotnie proporcjonalna do szerokości szczeliny: węższa szczelina powoduje większe rozmycie dyfrakcyjne. W siatce interferencja wielu szczelin zaostrza maksima, dlatego można rozróżniać bliskie sobie długości fal. Jeśli maksimum danego rzędu wypada w minimum obwiedni pojedynczej szczeliny, jest ono wygaszone. Przybliżenie Fraunhofera można uzyskać również za pomocą soczewki, umieszczając ekran w płaszczyźnie ogniskowej. Dyfrakcja ujawnia ograniczenia obrazowania i stanowi podstawę analizy widmowej.
\par

\subsection{27. Optyka geometryczna jako granica optyki falowej, przybliżenie paraxialne i lupa, teleskop, mikroskop}
\paragraph{Zakres.} Rozważa się przejście od optyki falowej do optyki geometrycznej, przybliżenie paraxialne oraz działanie lupy, teleskopu i mikroskopu. W granicy małej długości fali i dużych wymiarów układu, rozpraszanie dyfrakcyjne jest małe w porównaniu z rozmiarami układu, a promienie świetlne można traktować jako proste linie. Wtedy zamiast praw falowych stosuje się prawa odbicia i załamania; obraz tworzy się z prostoliniowych promieni. Przybliżenie paraxialne zakłada małe kąty względem osi optycznej i stosuje przybliżenie $\sin\theta\approx\tan\theta\approx\theta$.\par
Dla cienkiej soczewki obowiązuje równanie soczewkowe: $\frac{1}{f}=\frac{1}{x_o}+\frac{1}{x_i}$, gdzie $x_o$ jest odległością przedmiotu, $x_i$ obrazem, a $f$ ogniskową. Powiększenie liniowe wynosi $\beta=-x_i/x_o$. Lupa ma ogniskową małą i ustawiona jest tak, aby obraz powstał w odległości dobrego widzenia oka; powiększenie kątowe dla oka w odległości $D=25\,\mathrm{cm}$ wynosi $M\approx D/f$. Teleskop Keplera lub Galileusza ma układ obiektyw-okular; powiększenie kątowe to $M\approx -f_o/f_e$. Mikroskop złożony ma obiektyw o małej ogniskowej i okular; jego powiększenie całkowite wynosi $M\approx (L/f_o)(D/f_e)$.\par
To przybliżenie pozwala obliczać jasność, pole widzenia i rozdzielczość przy założeniu, że aberracje są pominięte. W rzeczywistych układach trzeba uwzględnić aberracje sferyczne i chromatyczne oraz ograniczenie dyfrakcyjne, które staje się ważne przy małych długościach fal i dużych aperturach. Tak więc optyka geometryczna jest granicą optyki falowej, bardzo przydatną do projektowania urządzeń optycznych.\par

Prawo Snella i równanie soczewki są przybliżeniami użytecznymi w układach, w których propagację można opisać promieniami, lecz nie przewidują wszystkich cech obrazu. Obraz rzeczywisty powstaje przez rzeczywiste przecięcie promieni, natomiast obraz pozorny przez przecięcie ich przedłużeń; to rozróżnienie pomaga opisywać konstrukcję obrazu. Teleskop zwiększa kąt, pod jakim widzimy odległy obiekt, ale nie musi zwiększać jego jasności powierzchniowej. Mikroskop uzyskuje duże powiększenie, gdy obiektyw tworzy powiększony obraz pośredni, a okular działa jak lupa. Rozdzielczość przyrządów nie rośnie dowolnie wraz z powiększeniem, ponieważ ograniczają ją dyfrakcja, apertura i długość fali. Przybliżenie paraxialne jest szczególnie dobre w pobliżu osi optycznej; dla promieni dalekich od osi pojawiają się aberracje. Dzięki temu optyka geometryczna daje prosty, ale świadomie ograniczony model instrumentów.
\par

\subsection{28. Odbicie i załamanie, całkowite wewnętrzne odbicie, dyspersja normalna i anormalna}
\paragraph{Zakres.} Rozpatruje się odbicie i załamanie światła na granicy ośrodków, warunek całkowitego wewnętrznego odbicia oraz zależność współczynnika załamania od częstotliwości. Prawo odbicia brzmi $\theta_r=\theta_i$, a prawo Snella ma postać $n_1\sin\theta_1=n_2\sin\theta_2$, $n=c/v$. Współczynnik załamania jest miarą zmiany szybkości propagacji fali w materiale. Jeśli światło przechodzi z ośrodka gęstszego do rzadszego, kąt załamania może przekroczyć $90^\circ$; całkowite wewnętrzne odbicie występuje przy $\sin\theta_c=n_2/n_1$ dla $n_1>n_2$. Zjawisko to leży u podstaw włókien optycznych i przekazów światła bez strat.\par
Dyspersja normalna oznacza, że współczynnik załamania rośnie wraz z częstotliwością, a anormalna — że maleje w pewnym zakresie. Dla wielu materiałów $n(\omega)$ zmienia się ze względu na odpowiedź atomową i polaryzację dipolową: w pobliżu rezonansu atomowego pojawiają się silne efekty zależne od częstotliwości. W pryzmacie odchylenie światła zależy od długości fali, więc światło białe rozszczepia się; w dyspersji normalnej czerwony promień odchyla się mniej niż fioletowy, a w dyspersji anormalnej zależność może się odwrócić.\par
Dyspersja ma kluczowe znaczenie w spektroskopii, soczewkach, włóknach i filtrach optycznych. Rozważania o współczynniku załamania w funkcji długości fali pozwalają korygować aberracje chromatyczne i projektować układy o wymaganych właściwościach. W optyce falowej jest to wynik zależności między polaryzacją materii a polem elektrycznym fali i ma ważne znaczenie w materiałoznawstwie.\par

Prawo odbicia wynika z ciągłości składowej stycznej wektora falowego na granicy ośrodków, a prawo załamania można uzasadnić niezmiennością częstości i zmianą prędkości fazowej. Przy kącie krytycznym fala załamana biegnie wzdłuż granicy; dla większego kąta nie ma propagującej fali załamanej w drugim ośrodku, choć w pobliżu granicy występuje pole zanikające. Włókno optyczne wykorzystuje całkowite wewnętrzne odbicie w rdzeniu otoczonym płaszczem o mniejszym współczynniku załamania. Dyspersja jest związana z odpowiedzią ośrodka na pole elektromagnetyczne; w pobliżu rezonansu następują także silne zmiany absorpcji. Określenie dyspersji anormalnej dotyczy lokalnej zależności współczynnika załamania od częstości i nie oznacza, że materiał nie może być przezroczysty w innych zakresach. Pryzmat rozszczepia wiązkę, ponieważ różne barwy mają różne prędkości fazowe w szkle, co jest ważne przy projektowaniu układów achromatycznych.
\par

\subsection{29. Polaryzacja światła, Malus, Brewster, dwójłomność, aktywność optyczna}
\paragraph{Zakres.} Rozpatruje się światło naturalne, spolaryzowane liniowo i kołowo, prawo Malusa, kąt Brewstera, zjawisko dwójłomności oraz aktywność optyczną materiałów. Światło naturalne po przejściu przez polaryzator ma natężenie $I=I_0/2$; po drugiej polaryzacji ustawionej pod kątem $\theta$ intensywność wynosi $I=I_0\cos^2\theta$. Jest to prawo Malusa. Gdy światło pada z jednego ośrodka na drugi, przy kącie Brewstera $\tan\theta_B=n_2/n_1$ fala odbita jest całkowicie spolaryzowana prostopadle do płaszczyzny padania. To zjawisko ma zastosowanie w polaryzatorach i analizatorach.\par
Dwójłomność występuje w kryształach anizotropowych, w których współczynnik załamania zależy od kierunku polaryzacji. Dla kryształu takich jak kalcyt wiązka świetlna rozszczepia się na dwa promienie o różnych prędkościach i różnych kierunkach. Dzięki temu można zbudować polaryzatory, kompensatory i modulatory światła. W materiałach aktywnych optycznie płaszczyzna polaryzacji światła jest obracana o kąt $\alpha$, zależny od grubości, stężenia i długości fali.\par
Aktywność optyczna jest zjawiskiem występującym w kryształach i roztworach niecentrosymetrycznych, w których odpowiedź dipolowa ma charakter wiązany. Dla substancji chiralnych istnieje dwójkierunkowa asymetria, co prowadzi do obrotu płaszczyzny polaryzacji. Zjawisko to ma zastosowanie w chemii, biologii i technologiach optycznych: do badania konfiguracji cząsteczek, kontroli jakości i badania związków organicznych.\par

Prawo Malusa opisuje projekcję wektora pola elektrycznego na oś transmisji analizatora; natężenie jest proporcjonalne do kwadratu amplitudy, stąd kwadrat cosinusa. Przy skrzyżowanych polaryzatorach idealnych światło nie przechodzi, ale w rzeczywistych materiałach pozostaje niewielka transmisja. Kąt Brewstera wynika z warunków brzegowych dla pola elektromagnetycznego i spełnia relację, że promień odbity i załamany są prostopadłe. W kryształach dwójłomnych promień zwyczajny i nadzwyczajny mają inne współczynniki załamania; ich prędkości i polaryzacje są różne. Aktywność optyczna może mieć różny znak dla enancjomerów tej samej substancji, co umożliwia określanie stężenia i składu mieszanin. Należy odróżniać obrót płaszczyzny polaryzacji od dwójłomności, która rozdziela polaryzacje i może wprowadzić między nimi przesunięcie fazy. Oba zjawiska wynikają z anizotropowej lub chiralnej odpowiedzi materiału.
\par

\subsection{30. Naturalna szerokość linii widmowej i czas życia, stany metastabilne}
\paragraph{Zakres.} Rozważa się naturalną szerokość linii spektralnych, związek z czasem życia stanów wzbudzonych oraz znaczenie stanów metastabilnych w laserach i spektroskopii. W mechanice kwantowej poziom energetyczny $E_2$ nie ma idealnie określonej energii, jeśli ma skończony czas życia $\tau$. Zasada nieoznaczoności daje $\Delta E\,\tau\gtrsim \hbar/2$, więc dla przejścia spontanicznego szerokość linii gaśnie z $\tau$. W idealnym przypadku naturalna szerokość linii Lorentza jest $\Delta\nu\approx 1/(2\pi\tau)$, a w praktyce dodatkowo rozmywają ją efekt Dopplera, ciśnienie i pole elektryczne.\par
Stan metastabilny ma bardzo długi czas życia w porównaniu ze zwykłymi stanami wzbudzonymi. W takich stanach elektrony mogą pozostawać wzbudzone przez dłuższy czas, co umożliwia gromadzenie inwersji ludności, niezbędnej do działania lasera. W laserze emitowany foton stymuluje kolejne przejścia z tego samego stanu wzbudzonego, co prowadzi do wzrostu amplitudy światła. To jest zasadnicza różnica między promieniowaniem spontanicznym i stymulowanym.\par
Stany metastabilne mają także zastosowanie w spektroskopii atomowej, gdzie słabe przejścia mogą być obserwowane na tle innych procesów. Długie czasy życia powodują, że linie są wąskie, a więc pomiar jest precyzyjny. Zależności między czasem życia, szerokością linii i prawem Einsteinów są fundamentem współczesnej fizyki atomowej i optyki kwantowej.\par

Szerokość naturalna jest nieodłączna od samego przejścia i nie da się jej usunąć przez doskonalenie aparatury, natomiast poszerzenie instrumentalne wynika z ograniczonej rozdzielczości przyrządu. Czas życia stanu jest średnim czasem przed emisją lub innym procesem zaniku, a nie gwarantowanym czasem, przez jaki pozostaje w nim każdy pojedynczy atom. Dla przejścia między dwoma stanami szerokość zależy od czasów życia obu poziomów, nie tylko poziomu górnego. Poszerzenie Dopplerowskie jest zwykle gaussowskie, a naturalne poszerzenie jest lorentzowskie; ich splot daje profil Voigta. Stan metastabilny bywa chroniony przez reguły wyboru, które sprawiają, że przejście jest mało prawdopodobne. Nie każdy długi czas życia oznacza automatycznie działanie laserowe: potrzebne są pompowanie, inwersja obsadzeń i rezonator optyczny. Analiza profilu linii pozwala rozdzielać mechanizmy poszerzenia i wyznaczać parametry atomów.
\par

\subsection{31. Rozpraszanie Rutherforda i wnioski o modelu atomu}
\paragraph{Zakres.} Rozważa się rozpraszanie cząstek alfa na cienkiej folii, rozkład kątowy i wnioski dotyczące rozmieszczenia ładunku w atomie. W doświadczeniu Rutherforda większość cząstek alfa przechodziła przez folię bez zakłóceń, ale niewielki odsetek ulegał rozproszeniu pod dużymi kątami, nawet powyżej $90^\circ$. Taki wynik był niezgodny z modelem Thomsona, w którym ładunek dodatni był rozproszony w całym atomie. Wniosek był prosty: dodatni ładunek i znaczna część masy są skoncentrowane w małym jądrze atomowym.\par
Rozpraszanie Rutherforda opisuje wzór
\[
\frac{d\sigma}{d\Omega}=\left(\frac{1}{4\pi\varepsilon_0}\right)^2\left(\frac{Zze^2}{4E_k}\right)^2\frac{1}{\sin^4(\theta/2)},
\]
który pokazuje, że rozpraszanie pod dużymi kątami jest rzadkie, ale możliwe. Gdy cząstka alfa zbliża się do jądra, siła Coulomba jest silna, a tor jest silnie odchylony. Miara rozkładu kątowego umożliwia oszacowanie rozmiaru jądra i liczby atomowej $Z$. Odkrycie to zmieniło obraz atomu: atom jest prawie pustą przestrzenią z elektronami poruszającymi się w dużej odległości od małego, gęstego jądra.\par
Model atomu Rutherforda był kluczowym krokiem do modelu Bohra i nowoczesnej mechaniki kwantowej. Pokazał, że najważniejsze zjawiska w atomach wynikają z centralnie lokalizowanego dodatniego ładunku i że elektron nie jest „wtopiony” w całą masę atomu. Wnioski z rozpraszania Rutherforda są nadal fundamentem fizyki jądrowej i atomowej.\par

Wzór Rutherforda zakłada klasyczne, sprężyste rozpraszanie na punktowym Coulombowskim jądrze, pojedyncze rozpraszanie oraz pomijalny wpływ elektronów w porównaniu z oddziaływaniem jądra. Kąt rozproszenia zależy od parametru zderzenia: im mniejsza odległość największego zbliżenia, tym silniejsze odchylenie. Natężenie dużych kątów jest małe, ponieważ tylko niewielka część cząstek alfa trafia w bliskie sąsiedztwo jądra. Zależność od energii cząstki pozwala sprawdzać, czy rozpraszanie jest zgodne z prawem Coulomba, a odchylenia od wzoru mogą ujawnić skończony rozmiar jądra. Doświadczenie nie wyjaśniało stabilności elektronów w atomie; klasyczny elektron krążący powinien promieniować i tracić energię. Ten problem wymagał późniejszych modeli kwantowych. Rozpraszanie Rutherforda dało zatem ważny obraz jądra, ale nie było kompletną teorią budowy atomu.
\par

\subsection{32. Rentgenowskie promieniowanie charakterystyczne i hamacyjne}
\paragraph{Zakres.} Rozważa się promieniowanie X powstające podczas bombardowania metalu elektronami oraz dwa główne typy widm: ciągłe promieniowanie hamowania i linie charakterystyczne. W przypadku zderzeń elektronów z materiałem, elektrony są gwałtownie hamowane w polu atomowym i emitują fotony o szerokim widmie. Granica krótkofalowa daje warunek $h\nu_{\max}=eU$, gdy elektron przyspieszony napięciem $U$ traci całą energię w jednym procesie. Zależność intensywności od długości fali jest ciągła, a „krótko falowy limit” wyraża się wzorem $\lambda_{\min}=\frac{hc}{eU}$.\par
Promieniowanie charakterystyczne pojawia się, gdy elektron z atomu docelowego wybije elektron z wewnętrznej powłoki, a zwolnione miejsce zapełnia elektron z wyższego poziomu. Różnica energii przejścia jest emitowana jako foton X. Tak powstają wyraźne linie widmowe, takie jak $K_\alpha$ i $K_\beta$, a ich energie zależą od pierwiastka. Jest to wykorzystywane w analizie chemicznej i identyfikacji materiałów. W spektroskopii rentgenowskiej każdy element ma charakterystyczny układ linii.\par
Promieniowanie X ma zastosowania w diagnostyce medycznej, tomografii komputerowej i rentgenografii kryształów. W biologii i medycynie jego pochłanianie zależnie od gęstości materiału tworzy kontrast obrazu; w kryształografii natomiast długość fali porównuje się z odległościami między płaszczyznami sieci. W ten sposób rentgenowskie promieniowanie X jest nieocenionym narzędziem badawczym w fizyce materii i medycynie.\par

Widmo hamowania ma postać ciągłą, ponieważ elektrony mogą oddawać różną część energii w wielu zderzeniach i przy różnych wartościach parametru zderzenia. Napięcie lampy rentgenowskiej określa graniczną energię fotonu, natomiast natężenie prądu wpływa przede wszystkim na liczbę elektronów i całkowitą intensywność promieniowania. Linie charakterystyczne pojawiają się dopiero po przekroczeniu energii wiązania danej powłoki; ich położenie zależy od atomu, a nie od napięcia lampy w szerokim zakresie powyżej progu. Po wybiciu elektronu powłokowego atom może również emitować elektron Augera zamiast fotonu rentgenowskiego. Widma X pozwalają odróżniać skład pierwiastkowy próbki i badać strukturę kryształów. W diagnostyce kontrast zależy od energii i składu tkanki, dlatego stosuje się filtry i właściwe napięcia, aby poprawić obraz oraz ograniczyć niepotrzebną dawkę.
\par

\subsection{33. Wzbudzanie i deekscytacja, przejścia gamma, konwersja wewnętrzna, Auger, Mössbauer}
\paragraph{Zakres.} Omawia się mechanizmy wzbudzania atomów i jąder, ich deekscytację oraz przejścia gamma i procesy towarzyszące. Wzbudzenie może nastąpić przez absorpcję fotonu, zderzenie z elektronami lub inne procesy. Stan wzbudzony ma wyższą energię niż stan podstawowy; deekscytacja może odbywać się przez emisję fotonu, przekazanie energii do otoczenia lub przejście między stanami jądrowymi. W atomach przejścia elektroniczne zwykle dają światło widzialne lub UV, natomiast dla jąder dominują promienie gamma.\par
Przejście gamma oznacza przejście między stanami jądra bez zmiany liczby nukleonów; energia przejścia jest emitowana w postaci fotonu gamma. Konwersja wewnętrzna polega na przekazaniu energii przejścia elektronowi z powłoki atomowej, który zostaje wyrzucony z atomu; Auger jest podobnym procesem, ale przy przekazywaniu energii innemu elektronowi. Oba zjawiska są konkurencyjne z emisją gamma i zależą od struktury elektronowej.\par
Zjawisko Mössbauera polega na bezretencjonalnej emisji lub absorpcji promieniowania gamma przez jądro osadzone w sieci krystalicznej. Dzięki bardzo wąskiej linii rezonansowej można mierzyć małe przesunięcia energii związane z polem elektrycznym i magnetycznym, jak również środowiskiem chemicznym. To pozwala na bardzo precyzyjne badania mechaniki kwantowej, magnetyzmu i struktury materiałów. Jest to ważne narzędzie zarówno w fizyce fundamentalnej, jak i w materiałoznawstwie.\par

Wzbudzenie atomu i wzbudzenie jądra dotyczą różnych skal energii, dlatego charakterystyczne energie fotonów elektronowych i gamma różnią się zwykle o wiele rzędów wielkości. Konwersja wewnętrzna nie jest emisją fotonu: elektron otrzymuje energię przejścia pomniejszoną o energię wiązania swojej powłoki, a powstała wakancja może następnie wywołać promieniowanie X lub elektrony Augera. W procesie Augera energia relaksacji atomu przekazuje się kolejnemu elektronowi, który opuszcza atom; proces ten konkuruje z emisją charakterystycznego promieniowania X. W efekcie po przejściu jądrowym mogą powstawać linie elektronów i fotonów, użyteczne w analizie materiałów. Efekt Mössbauera jest możliwy tylko dla wybranych izotopów i przejść gamma, gdy odrzut pochłania sieć jako całość, a nie pojedyncze jądro. Dzięki temu energia fotonu może pokryć się z energią rezonansu absorpcji. Zmiany przesunięcia izomerycznego ujawniają gęstość elektronową przy jądrze.
\par

\subsection{34. Odkrycie jądra atomowego, struktura i własności jądra}
\paragraph{Zakres.} Rozważa się historię odkrycia jądra atomowego, strukturę jądra oraz jego podstawowe właściwości: ładunek, rozmiar, masa, spin i gęstość. Doświadczenie Rutherforda wykazało, że atom ma małe, gęste, dodatnio naładowane jądro, wokół którego krążą elektrony. Rozmiar jądra ma rząd $10^{-15}$–$10^{-14}\,\mathrm{m}$, czyli jest 10 000 do 100 000 razy mniejszy niż rozmiar atomu. Znaczna część masy atomu znajduje się w jądrze, a elektrony zajmują obszar zewnętrzny.\par
Jądro składa się z protonów i neutronów, zwanych nukleonami. Proton ma ładunek $+e$, neutron jest elektrycznie obojętny. Liczba protonów $Z$ określa ładunek i numer atomowy pierwiastka, a liczba neutronów $N$ jest istotna dla izotopów. Liczba masowa $A=Z+N$. Jądra mają spin, moment magnetyczny, poziomy energetyczne i określoną energię wiązania. Gęstość materii jądrowej jest ogromna, co wynika z silnego oddziaływania jądrowego i bardzo małych rozmiarów.\par
Struktura jądra jest opisana przez mechanikę kwantową; nukleony zajmują stany o określonych energiach, podobnie jak elektrony w atomie. Właściwości jąder są badane metodami spektroskopii jądrowej, rozpraszania i reakcji jądrowych. Zrozumienie budowy jądra miało kluczowy wpływ na rozwój fizyki atomowej, chemii, energii jądrowej i medycyny nuklearnej.\par

Promień jądra rośnie w przybliżeniu proporcjonalnie do pierwiastka sześciennego z liczby masowej, co wskazuje na prawie stałą gęstość materii jądrowej. Energia wiązania to energia potrzebna do rozdzielenia jądra na swobodne nukleony; deficyt masy odpowiada tej energii przez zależność równoważności masy i energii. Izotopy tego samego pierwiastka mają jednakowy ładunek, lecz różną liczbę neutronów, dlatego mogą różnić się stabilnością i własnościami jądrowymi. Oddziaływanie silne wiąże protony i neutrony na krótkich odległościach, pokonując odpychanie elektrostatyczne protonów. Jądro nie jest małym skupiskiem protonów bez neutronów: neutrony zwiększają wiązanie bez dokładania odpychania Coulombowskiego. Własności takie jak spin i moment magnetyczny pomagają badać strukturę poziomów, a rozpraszanie i reakcje jądrowe dostarczają informacji o rozmiarze i rozkładzie gęstości ładunku.
\par

\subsection{35. Modele jądra: kroplowy i powłokowy}
\paragraph{Zakres.} Przedstawia się dwa klasyczne modele jąder: model kroplowy i model powłokowy, ich założenia, sukcesy i ograniczenia. Model kroplowy traktuje jądro jako niemal nieściśliwą kroplę płynu, w której energia wiązania składa się z członów objętościowych, powierzchniowych, Coulombowskich i parzystości: $E_B(A,Z)=a_VA-a_SA^{2/3}-a_C\frac{Z^2}{A^{1/3}}-a_A\frac{(N-Z)^2}{A}+\delta(A,Z)$. Ten opis dobrze oddaje ogólną zależność energii wiązania od liczby nukleonów oraz zjawiska rozszczepienia ciężkich jąder.\par
Model powłokowy zakłada, że nukleony zajmują kwantowe poziomy energetyczne, podobnie jak elektrony w atomie. Pełne powłoki odpowiadają szczególnie stabilnym jądom, określanym jako liczby magiczne: 2, 8, 20, 28, 50, 82, 126. Modele powłokowe wyjaśniają stabilność izotopów, poziomy energetyczne jąder, momenty magnetyczne i promieniowanie gamma. Przykładowo, jądra z liczbami magicznymi są bardziej trwałe niż sąsiednie nuklidy.\par
Każdy z modeli ma ograniczenia: kroplowy dobrze opisuje ogólną energię i rozszczepienie, ale nie wyjaśnia struktury poziomów i liczb magicznych; model powłokowy jest lepszy dla pojedynczych stanów i własności strukturalnych, ale mniej wygodny dla bardzo ciężkich jąder. Dlatego w praktyce wykorzystuje się oba modele jako komplementarne opisowe.\par

Model kroplowy prowadzi do półempirycznego wzoru masowego, którego składniki mają proste interpretacje: objętość wiąże nukleony z sąsiadami, powierzchnia zmniejsza liczbę sąsiadów, odpychanie Coulombowskie osłabia wiązanie protonów, a człon asymetrii odzwierciedla zasadę Pauliego. Poprawka parzystości wiąże się z tendencją do parowania protonów i neutronów. Model kroplowy przewiduje również, że odkształcenie jądra może prowadzić do rozszczepienia, jeśli zysk energii powierzchniowej przeważy nad wzrostem energii elektrostatycznej. W modelu powłokowym konieczne jest uwzględnienie silnego sprzężenia spin-orbita, aby odtworzyć obserwowane liczby magiczne. Modele nie są sprzecznymi konkurentami: rzeczywiste jądra mają zarówno kolektywne ruchy podobne do kropli, jak i kwantową strukturę pojedynczych nukleonów. W wielu zastosowaniach łączy się oba opisy, co pozwala interpretować deformacje i widma.
\par

\subsection{36. Rozpad promieniotwórczy: rodzaje i teoria}
\paragraph{Zakres.} Rozważa się aktywność promieniotwórczą, rodzaje rozpadów oraz prawo rozpadu i stałą rozpadu. Rozpad promieniotwórczy jest procesem samorzutnym, w którym niestabilne jądro przechodzi w stan o niższej energii, emitując cząstkę lub foton. Najważniejsze typy to rozpad alfa ($^4\mathrm{He}$), beta minus ($n\to p+e^-+\bar\nu$), beta plus ($p\to n+e^++\nu$), wychwyt elektronu i emisja gamma. Każdy rozpad zachowuje ładunek, liczbę nukleonów i pęd, choć w rozpadach beta zachodzi także zachowanie leptonowe i udział neutrin.\par
Prawo rozpadu ma postać $N(t)=N_0e^{-\lambda t}$, $T_{1/2}=\ln2/\lambda$, $A=\lambda N$, gdzie $A$ to aktywność, $\lambda$ stała rozpadu, a $T_{1/2}$ czas połowicznego zaniku. Rozpad jest zjawiskiem statystycznym: dla pojedynczego jądra nie da się przewidzieć chwili rozpadu, ale dla dużej liczby jąder opis statystyczny jest całkowicie poprawny. Typowy izotop ma charakterystyczny czas połowicznego zaniku.\par
Rozpad alfa zmniejsza liczbę masową o 4 i ładunek o 2; beta minus zamienia neutron w proton, beta plus — proton w neutron; gamma zachodzi między stanami jądrowymi bez zmiany liczby nukleonów. Zjawiska te mają kluczowe znaczenie w datowaniu radiowęglowym, ochronie radiologicznej, diagnostyce medycznej i produkcji energii jądrowej.\par

Prawo wykładnicze wynika z założenia, że prawdopodobieństwo rozpadu każdego nieuległego jeszcze rozpadowi jądra w jednostce czasu jest stałe i niezależne od jego wieku. Aktywność mierzy się w bekerelach, przy czym jeden bekerel oznacza jeden rozpad na sekundę; nie jest to bezpośrednio miara dawki pochłoniętej przez człowieka. W rzeczywistej próbce mogą powstawać produkty rozpadu, które same są promieniotwórcze, tworząc szeregi i stan równowagi między nuklidami. Rozpad alfa zachodzi często w ciężkich jądrach, a mechanika kwantowa wyjaśnia jego prawdopodobieństwo tunelowaniem przez barierę Coulomba. Widmo energii cząstek beta jest ciągłe, ponieważ energia dzieli się między elektron lub pozyton i neutrino. Rozpad gamma zwykle następuje po innym rozpadzie, gdy jądro potomne pozostaje wzbudzone. Te różnice pozwalają identyfikować kanał przemiany i wyznaczać skład izotopowy.
\par

\subsection{37. Oddziaływanie promieniowania z materią}
\paragraph{Zakres.} Rozważa się mechanizmy pochłaniania promieniowania przez materię, zasięg cząstek jonizujących i niejonizujących oraz skutki biologiczne. Promieniowanie alfa ma dużą moc jonizacji, ale ma krótki zasięg; beta jest bardziej przenikliwe; gamma i promienie X przechodzą przez materię głębiej. Prawo osłabiania ma postać $I=I_0e^{-\mu x}$, gdzie $\mu$ jest współczynnikiem osłabienia liniowego. Głównymi mechanizmami są fotoelektryczny efekt, rozpraszanie Comptona i tworzenie par elektron-pozyton.\par
W niskich energiach dominuje efekt fotoelektryczny, przy średnich — rozpraszanie Comptona, a przy energiach powyżej $1.022\,\mathrm{MeV}$ — produkcja par. Warto zauważyć, że różne typy promieniowania mają różne mechanizmy interakcji z materią i więc różny zasięg. Dla gatunku biologicznego decyduje to o skuteczności ochrony, np. o konieczności osłon wykonanych z ołów, betonu lub materiałów lekkich, zależnie od rodzaju promieniowania.\par
Skutki biologiczne wynikają z jonizacji i tworzenia wolnych rodników, co może uszkadzać DNA i tkanki. Ochrona przed promieniowaniem opiera się na odległości od źródła, ograniczeniu czasu ekspozycji i ekranowaniu. Zjawiska te są również wykorzystywane w diagnostyce, terapii onkologicznej i spektroskopii. Oddziaływanie promieniowania z materią jest zatem fundamentem zarówno fizyki jądrowej, jak i medycyny.\par

Dla cząstek naładowanych energia jest przekazywana głównie przez zderzenia z elektronami atomowymi, powodując wzbudzenia i jonizację; ciężkie cząstki, takie jak alfa, tracą energię intensywnie na krótkiej drodze. Elektrony mogą dodatkowo emitować promieniowanie hamowania, szczególnie w materiałach o dużej liczbie atomowej. Fotony gamma nie mają określonego końcowego zasięgu w taki sam sposób jak cząstki naładowane; prawdopodobieństwo oddziaływania opisuje współczynnik osłabienia i zależy od energii oraz składu ośrodka. Neutrony oddziałują przede wszystkim z jądrami i mogą być spowalniane przez zderzenia, a następnie pochłaniane. W ochronie dobiera się osłony do rodzaju promieniowania: alfa zatrzymuje papier lub warstwa zewnętrzna skóry, beta wymaga lekkich materiałów, a gamma — gęstych osłon. Pochłonięta dawka i biologiczna skuteczność to różne wielkości, dlatego bezpieczeństwo ocenia się według dawki i rodzaju promieniowania.
\par

\clearpage
\section{Fizyka teoretyczna}
\subsection{1. Formalizm Lagrange’a i równania Eulera--Lagrange’a: oscylator harmoniczny}
Zakres: opis ruchu za pomocą współrzędnych i działania, bez bezpośredniego używania sił Newtona. Dla układu o współrzędnych uogólnionych \(q_i\) i prędkościach \(\dot q_i\) lagranżjan definiujemy jako \(L=T-V\), gdzie \(T\) jest energią kinetyczną, a \(V\) potencjalną. Rzeczywista trajektoria w przedziale czasu \([t_1,t_2]\) czyni działanie \(S=\int_{t_1}^{t_2}L\,dt\) stacjonarnym przy wariacjach trajektorii zerujących się na końcach. Z warunku \(\delta S=0\), po całkowaniu przez części wyrazu z \(\delta\dot q_i\), otrzymujemy równania Eulera--Lagrange’a:
\[
\frac{d}{dt}\frac{\partial L}{\partial \dot q_i}-\frac{\partial L}{\partial q_i}=0.
\]
Dla jednowymiarowego oscylatora o masie \(m\), stałej sprężystości \(k>0\) i wychyleniu \(x\) przyjmujemy \(T=m\dot x^2/2\), \(V=kx^2/2\). Zatem
\[
L=\frac{m\dot x^2}{2}-\frac{kx^2}{2},\qquad m\ddot x+kx=0.
\]
Po wprowadzeniu \(\omega=\sqrt{k/m}\) rozwiązanie ma postać \(x(t)=A\cos(\omega t)+B\sin(\omega t)\); stałe wyznaczają warunki początkowe. Energia \(E=m\dot x^2/2+kx^2/2\) jest zachowana, bo \(L\) nie zależy jawnie od czasu. Formalizm zakłada układ klasyczny i gładki lagranżjan. Jest szczególnie użyteczny, gdy więzy upraszczają wybór współrzędnych lub symetria prowadzi do zachowania pędu uogólnionego.

Przykład pokazuje, jak zasada wariacyjna odtwarza dynamikę znaną z drugiej zasady Newtona, ale jej treść jest szersza niż samo przepisanie siły. Dodanie do lagranżjanu pochodnej zupełnej po czasie nie zmienia równań ruchu, ponieważ zmienia działanie jedynie o wartości na końcach przedziału. Można więc wybierać równoważne postacie opisu, nie zmieniając przewidywanej trajektorii. Jeśli współrzędna nie występuje jawnie w lagranżjanie, jej pęd uogólniony jest zachowany; dla oscylatora położenie nie jest cykliczne, więc pęd \(m\dot x\) nie pozostaje stały. Przy \(x(0)=x_0\) i \(\dot x(0)=0\) ciało rozpoczyna ruch w skrajnym położeniu, a następnie przechodzi przez równowagę z maksymalną szybkością. W równowadze sprężyna nie wykonuje pracy wypadkowej w całym okresie, choć energia kinetyczna i potencjalna cyklicznie zamieniają się. Ten przykład służy też jako lokalny model stabilnego ruchu w pobliżu minimum potencjału: rozwinięcie dowolnego gładkiego potencjału do drugiego rzędu daje w przybliżeniu oscylator harmoniczny.
\subsection{2. Stopnie swobody i współrzędne uogólnione: cząstka swobodna na sferze}
Zakres: liczenie niezależnych parametrów ruchu i wybór współrzędnych dopasowanych do więzów. W przestrzeni trójwymiarowej punkt materialny ma trzy stopnie swobody. Każdy niezależny więz holonomiczny, czyli równanie \(f(\mathbf r,t)=0\), usuwa jeden stopień swobody; cząstka ograniczona do powierzchni sfery o promieniu \(R\) spełnia \(x^2+y^2+z^2=R^2\), więc ma dwa. Naturalnymi współrzędnymi uogólnionymi są kąty biegunowe \(\theta,\phi\), a nie kartezjańskie \(x,y,z\). Zakładamy stały promień, brak potencjału i masę \(m\). Parametryzacja
\[
\mathbf r=R(\sin\theta\cos\phi,\sin\theta\sin\phi,\cos\theta)
\]
daje \(v^2=R^2(\dot\theta^2+\sin^2\theta\,\dot\phi^2)\), zatem \(L=mR^2(\dot\theta^2+\sin^2\theta\,\dot\phi^2)/2\). Równanie dla \(\theta\) oraz dla \(\phi\) wynika z równań Eulera--Lagrange’a. Ponieważ \(\phi\) nie występuje jawnie w \(L\), jest współrzędną cykliczną i
\[
p_\phi=\frac{\partial L}{\partial\dot\phi}=mR^2\sin^2\theta\,\dot\phi
\]
jest stałe. To składowa momentu pędu względem osi \(z\). Równanie dla \(\theta\) brzmi \(\ddot\theta-\sin\theta\cos\theta\,\dot\phi^2=0\). Ruch pozostaje po wielkim kole (przy stałym wektorze momentu pędu), co można również uzyskać z równania siły reakcji więzu. Współrzędne kątowe są osobliwe na biegunach, ale nie oznacza to fizycznej osobliwości ruchu.

Współrzędne uogólnione nie muszą mieć wymiaru długości: kąt jest bezwymiarowy, a jego prędkość wyraża się w radianach na sekundę. Liczba stopni swobody oznacza wymiar przestrzeni konfiguracji po uwzględnieniu więzów, a nie liczbę składowych wektora położenia w wybranym układzie. Na sferze dwa kąty określają punkt poza biegunami, gdzie współrzędna azymutalna traci jednoznaczność; do opisu okolicy bieguna można wybrać inną mapę współrzędnych. Osobliwość tej parametryzacji nie tworzy ani nie usuwa fizycznego stopnia swobody. Cząstka jest swobodna w tym sensie, że nie działa na nią potencjał wzdłuż powierzchni, ale więz sferyczny wymaga siły reakcji normalnej do powierzchni. Ta siła nie wykonuje pracy, ponieważ prędkość jest styczna do sfery. Zachowany pęd \(p_\phi\) wynika z symetrii obrotowej wokół osi \(z\); odpowiadająca jej symetria przestrzeni daje zachowanie odpowiedniej składowej momentu pędu. Dobór współrzędnych uwzględniających więz pozwala nie rozwiązywać osobno siły reakcji, co upraszcza zarówno zapis energii, jak i równania ruchu.
\subsection{3. Hamiltonian i równania Hamiltona: rozwiązanie oscylatora}
Zakres: zamiana opisu \((q,\dot q)\) na opis położenia i pędu kanonicznego. Pęd kanoniczny definiuje się przez \(p_i=\partial L/\partial\dot q_i\), a transformację Legendre’a przez \(H=\sum_i p_i\dot q_i-L\), przy założeniu, że zależność pędu od prędkości można odwrócić. Hamiltonian jest funkcją \(q,p,t\), a jego równania ruchu to
\[
\dot q_i=\frac{\partial H}{\partial p_i},\qquad
\dot p_i=-\frac{\partial H}{\partial q_i}.
\]
Dla oscylatora z \(L=m\dot x^2/2-kx^2/2\) mamy \(p=m\dot x\), a więc \(H=p^2/(2m)+kx^2/2\). W tym przykładzie Hamiltonian jest energią całkowitą. Równania Hamiltona dają \(\dot x=p/m\) oraz \(\dot p=-kx\). Różniczkując pierwsze z nich i podstawiając drugie, otrzymujemy
\[
\ddot x+\omega^2x=0,\qquad \omega=\sqrt{k/m}.
\]
Rozwiązując równanie drugiego rzędu i używając \(p=m\dot x\), przy warunkach \(x(0)=x_0,\ p(0)=p_0\) dostajemy
\[
x(t)=x_0\cos\omega t+\frac{p_0}{m\omega}\sin\omega t,\quad
p(t)=p_0\cos\omega t-m\omega x_0\sin\omega t.
\]
Trajektoria w przestrzeni fazowej \((x,p)\) jest elipsą stałego \(H\); jej kształt odzwierciedla różne jednostki współrzędnej i pędu. Jeżeli \(H\) nie zależy jawnie od czasu, to \(dH/dt=0\). Hamiltonian bywa równy energii, lecz w układach zależnych od czasu albo przy szczególnych więzach nie należy utożsamiać tych pojęć bez sprawdzenia.

Rozwiązanie jawne pozwala sprawdzić, że warunki początkowe są spełnione: dla \(t=0\) otrzymujemy \(x_0\) oraz \(p_0\). W chwili ćwierci okresu początkowe wychylenie zamienia się w pęd, co ilustruje obrót elipsy fazowej. W punkcie równowagi potencjał jest najmniejszy, ale prędkość może być największa; w skrajnym wychyleniu pęd chwilowo znika. Transformacja Legendre’a wymaga odwracalnej zależności \(p=m\dot x\), co dla dodatniej masy jest spełnione. W opisie Hamiltona położenie i pęd są niezależnymi zmiennymi stanu, a równania są pierwszego rzędu, dlatego do wyznaczenia ewolucji podaje się obie wartości początkowe. Dla dowolnej funkcji obserwabli \(F(x,p,t)\) jej pochodną wzdłuż ruchu można obliczyć z równania Poissona i jawnej zależności od czasu. W szczególności stały Hamiltonian jest zachowany, ponieważ nie ma jawnej zależności czasowej. Przybliżenie harmoniczne jest ważne dla potencjału sprężystego liniowego; w układach nieliniowych zależność między okresem a energią może być inna.
\subsection{4. Wahadło: współrzędne i pęd uogólniony, równania Hamiltona, małe drgania}
Zakres: wahadło matematyczne o masie punktowej \(m\), nieważkiej nici długości \(\ell\), stałym punkcie zawieszenia i jednorodnym polu grawitacyjnym \(g\). Jedyny stopień swobody opisuje kąt \(\theta\) od pionu skierowanego w dół. Pomijamy opór i zakładamy nierozciągliwą nić. Prędkość masy ma wartość \(\ell\dot\theta\), a po wyborze zera energii potencjalnej w najniższym położeniu
\[
T=\frac12m\ell^2\dot\theta^2,\quad
V=mg\ell(1-\cos\theta),\quad L=T-V.
\]
Pęd uogólniony i Hamiltonian wynoszą
\[
p_\theta=\frac{\partial L}{\partial\dot\theta}=m\ell^2\dot\theta,\qquad
H=\frac{p_\theta^2}{2m\ell^2}+mg\ell(1-\cos\theta).
\]
Równania Hamiltona to \(\dot\theta=p_\theta/(m\ell^2)\) oraz \(\dot p_\theta=-mg\ell\sin\theta\). Łącząc je, otrzymujemy \(\ddot\theta+(g/\ell)\sin\theta=0\), równoważne równaniu Lagrange’a. Dla małych wychyleń w radianach \(|\theta|\ll1\) stosujemy \(\sin\theta\simeq\theta\), \(1-\cos\theta\simeq\theta^2/2\). Wtedy
\[
\ddot\theta+\omega_0^2\theta=0,\qquad \omega_0=\sqrt{g/\ell},\quad
T_0=2\pi\sqrt{\ell/g}.
\]
Okres małych drgań nie zależy w tym przybliżeniu od masy ani amplitudy. Przy większych wychyleniach przybliżenie zawodzi, okres rośnie z amplitudą, a pełne rozwiązanie wymaga całki eliptycznej. Pęd ma wymiar momentu pędu, bo współrzędna \(\theta\) jest bezwymiarowa. Zachowanie \(H\) wynika z braku jawnej zależności od czasu.

Energia potencjalna wahadła rośnie przy unoszeniu masy, a jej najniższe położenie jest stabilne, gdyż małe wychylenie wywołuje moment siły przywracający ciało do pionu. Równanie dokładne zawiera sinus, dlatego zasada superpozycji rozwiązań nie obowiązuje przy dużych kątach. W granicy małych drgań sinus można zastąpić kątem wyrażonym w radianach; rozwiązanie jest wtedy sinusoidalne, a amplituda zależy od początkowego kąta i prędkości kątowej. Jeśli puszczamy wahadło bez prędkości z kąta \(\theta_{\max}\), energia początkowa jest potencjalna, a przy przejściu przez pion cała dostępna energia staje się kinetyczna. Naprężenie nici nie jest dodatkowym stopniem swobody, ponieważ długość pozostaje stała i więz geometryczny już uwzględnia jego skutek. W eksperymencie rzeczywiste wahadło ma skończony rozmiar, opór powietrza i straty w zawieszeniu; wówczas okres może odbiegać od idealnej wartości, a amplituda z czasem maleje. Dla małych amplitud wzór na okres umożliwia wyznaczenie lokalnego przyspieszenia grawitacyjnego z pomiaru długości i okresu. To zastosowanie wymaga jednak precyzyjnego pomiaru długości od osi obrotu do środka masy.
\subsection{5. Ruch ładunku próbnego w statycznych polach \(\mathbf E,\mathbf B\): siła Lorentza}
Zakres: klasyczny ruch punktowego ładunku \(q\) o masie \(m\) w zadanych, niezależnych od czasu polach elektrycznym i magnetycznym. Ładunek próbny oznacza, że jego własny wpływ na źródła pól pomijamy. Równanie Lorentza ma postać
\[
\mathbf F=q(\mathbf E+\mathbf v\times\mathbf B),\qquad
m\frac{d\mathbf v}{dt}=q(\mathbf E+\mathbf v\times\mathbf B).
\]
Pole elektryczne może zmieniać energię kinetyczną, ponieważ moc wynosi \(q\mathbf E\cdot\mathbf v\). Część magnetyczna jest prostopadła do prędkości i nie wykonuje pracy: \(\mathbf v\cdot(\mathbf v\times\mathbf B)=0\). Dla samego jednorodnego pola \(\mathbf B=B\hat{\mathbf z}\), prędkość równoległa do pola pozostaje stała, a składowa prostopadła wykonuje ruch kołowy o częstości cyklotronowej \(\omega_c=|q|B/m\) i promieniu \(r_L=m v_\perp/(|q|B)\). Cały tor jest helisą; znak \(q\) ustala kierunek obrotu. Dla samego jednorodnego pola \(\mathbf E\) przyspieszenie wynosi \(q\mathbf E/m\).

Przy jednorodnych, prostopadłych polach \(\mathbf E\perp\mathbf B\) występuje dryf guiding-center o prędkości
\[
\mathbf v_D=\frac{\mathbf E\times\mathbf B}{B^2},
\]
niezależnej od znaku i wartości ładunku w nierelatywistycznym przybliżeniu. Wynika to po przejściu do układu poruszającego się z tą prędkością, w którym pole elektryczne znika. Przybliżenie klasyczne jest dobre, gdy efekty kwantowe i relatywistyczne są zaniedbywalne, a pola można traktować jako zewnętrzne. W polu niejednorodnym ruch może łączyć lokalną spiralę z dryfami; proste wzory na promień i dryf wymagają lokalnie stałego pola.

Zasada superpozycji pozwala rozdzielić ruch na składowe, ale w obecności obu pól trzeba rozwiązać równanie wektorowe, bo kierunek siły magnetycznej zmienia się wraz z prędkością. Dla prędkości początkowej równoległej do jednorodnego pola magnetycznego nie ma siły magnetycznej, a cząstka porusza się ruchem jednostajnym wzdłuż linii pola. Składowa prostopadła powoduje zakrzywienie toru, nie zmianę energii; zatem sam magnes może zmieniać kierunek pędu, lecz nie wartość energii kinetycznej. W praktyce separatory mas wykorzystują zależność promienia krzywizny od stosunku pędu do ładunku, a spektrometry mierzące energię wykorzystują pole elektryczne. Wielkość \(q\mathbf E\) jest siłą, natomiast \(q\mathbf v\times\mathbf B\) ma te same jednostki dzięki jednostkom indukcji magnetycznej. Wzór na dryf jest istotny dla naładowanych cząstek w plazmie, lecz zakłada stałe pola i pomija efekty graniczne oraz bezwładność przy gwałtownej zmianie warunków. Dla cząstek relatywistycznych trzeba zastąpić \(m\mathbf v\) pędem relatywistycznym, chociaż siła Lorentza zachowuje tę samą postać.
\subsection{6. Prawo Gaussa dla \(\mathbf E\) i \(\mathbf B\): postać całkowa, różniczkowa i prawo Coulomba}
Zakres: związki strumieni pól z ładunkiem oraz brak obserwowanych monopoli magnetycznych w klasycznej elektrodynamice. W próżni prawo Gaussa dla elektryczności i magnetyzmu ma postać
\[
\oint_{\partial V}\mathbf E\cdot d\mathbf A=\frac{Q_{\mathrm{wew}}}{\varepsilon_0},
\qquad \oint_{\partial V}\mathbf B\cdot d\mathbf A=0.
\]
Pierwsze mówi, że całkowity strumień \(\mathbf E\) przez zamkniętą powierzchnię zależy od ładunku wewnętrznego, a nie od jego położenia w tej powierzchni; drugie, że linie \(\mathbf B\) nie mają źródeł ani ujść. Z twierdzenia Gaussa w rachunku wektorowym wynika dla dowolnej objętości
\[
\nabla\cdot\mathbf E=\rho/\varepsilon_0,\qquad \nabla\cdot\mathbf B=0,
\]
gdzie \(\rho\) jest gęstością ładunku. Dla punktowego ładunku \(q\) symetria sferyczna wymusza radialne pole o stałym natężeniu na sferze promienia \(r\). Podstawienie do prawa całkowego daje \(E(4\pi r^2)=q/\varepsilon_0\), zatem
\[
\mathbf E(\mathbf r)=\frac{q}{4\pi\varepsilon_0r^2}\hat{\mathbf r}.
\]
To pole Coulomba; dla wielu ładunków stosuje się superpozycję. W próżni \(\varepsilon_0\) oznacza przenikalność elektryczną, \(d\mathbf A\) jest skierowanym elementem powierzchni, a \(Q_{\rm wew}=\int_V\rho\,dV\). Prawo Coulomba wynika tu z prawa Gaussa dopiero po przyjęciu symetrii sferycznej (i stacjonarności), a nie samo z siebie dla dowolnego rozkładu ładunku. W materii należy odpowiednio rozróżnić ładunki swobodne i związane.

W zastosowaniach prawo Gaussa najłatwiej wykorzystać wtedy, gdy rozkład źródeł ma symetrię sferyczną, cylindryczną lub płaską; wtedy pole ma stałą wartość na odpowiednio dobranej powierzchni. Dla niesferycznego układu ładunków prawo pozostaje prawdziwe, lecz sam strumień nie wystarcza do wyznaczenia lokalnej wartości pola bez dodatkowych informacji. Ładunki zewnętrzne mogą zmieniać pole na powierzchni, ale ich wkład do całkowitego strumienia przez zamkniętą powierzchnię znosi się. W próżni linie pola elektrycznego zaczynają się na ładunkach dodatnich i kończą na ujemnych lub w nieskończoności; linie pola magnetycznego tworzą zamknięte pętle. Drugie równanie jest empirycznym stwierdzeniem braku monopoli magnetycznych w klasycznym elektromagnetyzmie, a nie konsekwencją symetrii sferycznej. W ośrodku liniowym i jednorodnym można zamiast pola \(\mathbf E\) używać indukcji elektrycznej \(\mathbf D\), przy czym przenikalność materiału modyfikuje związek między polem i źródłami swobodnymi. Prawo Coulomba dla punktowych ładunków prowadzi po pomnożeniu pola przez ładunek próbny do prawa siły, a kierunek zależy od znaków obu ładunków.
\subsection{7. Prawo Faradaya: postać całkowa i różniczkowa, SEM indukcji w obwodzie}
Zakres: indukcja elektromagnetyczna w klasycznej elektrodynamice. Zmiana strumienia magnetycznego przez powierzchnię ograniczoną obwodem powoduje elektromotoryczną siłę indukcji (SEM). Dla nieruchomej pętli \(C=\partial S\) prawo Faradaya--Lenza zapisujemy
\[
\mathcal E=\oint_C\mathbf E\cdot d\boldsymbol\ell
=-\frac{d}{dt}\int_S\mathbf B\cdot d\mathbf A
=-\frac{d\Phi_B}{dt}.
\]
Znak minus wyraża regułę Lenza: indukowany efekt przeciwdziała zmianie strumienia, która go wywołała. Stosując twierdzenie Stokesa dla ustalonej powierzchni, otrzymujemy
\[
\nabla\times\mathbf E=-\frac{\partial\mathbf B}{\partial t}.
\]
Pole elektryczne wywołane zmiennym polem magnetycznym jest wirowe i nie musi być zachowawcze; jego całka po zamkniętym obwodzie może być niezerowa. Dla cewki o \(N\) zwojach, których strumienie są jednakowe, \(\mathcal E=-N\,d\Phi_B/dt\). W obwodzie o oporze \(R\), przy założeniu quasi-stacjonarnym i pominięciu samoindukcji, prąd ma wartość \(I=\mathcal E/R\). Przykładowo obracająca się ramka o polu \(A\) w jednorodnym polu \(B\) ma \(\Phi_B=BA\cos(\omega t)\), więc wytwarza sinusoidalną SEM o amplitudzie \(BA\omega\). Jeśli obwód się porusza, ogólna SEM uwzględnia również siłę magnetyczną na poruszające się ładunki; uproszczona postać strumieniowa wymaga poprawnego uwzględnienia ruchu granicy. Wzór różniczkowy jest lokalny i obowiązuje punktowo dla gładkich pól.

Zmianę strumienia można uzyskać przez zmianę wartości pola, powierzchni pętli albo kąta między polem i jej normalną. Każdy z tych sposobów prowadzi do indukcji, o ile strumień rzeczywiście zmienia się w czasie. SEM jest pracą sił elektromagnetycznych przypadającą na jednostkowy ładunek wzdłuż obwodu, a nie różnicą potencjału elektrostatycznego między dwoma punktami. W statycznym polu elektrostatycznym całka po pętli znika, natomiast wir pola elektrycznego wywołanego zmianą magnetyczną może być niezerowy. Reguła Lenza wynika z zachowania energii: prąd indukowany nie może samorzutnie wzmacniać zmiany, która go wywołała, bez dopływu pracy z zewnątrz. W generatorze obrót ramki wymaga dostarczenia pracy mechanicznej, która zamienia się w energię elektryczną i ciepło w obciążeniu. Przy wielu uzwojeniach należy ustalić orientację zwojów; odwrócenie normalnej powierzchni odwraca jednocześnie znak strumienia i kierunek obiegu, nie zmieniając fizycznego przewidywania. Dla prądu zmiennego istotna staje się indukcyjność własna, więc rzeczywisty prąd nie musi być po prostu równy chwilowej SEM podzielonej przez opór.
\subsection{8. Prawo Ampère’a--Maxwella i ładowanie kondensatora płaskiego}
Zakres: pole magnetyczne wytwarzane zarówno przez prąd przewodzenia, jak i zmienne pole elektryczne. Prawo Ampère’a dla prądów stacjonarnych wymaga poprawki Maxwella, aby zgadzało się z zachowaniem ładunku. W próżni postać całkowa i różniczkowa brzmią
\[
\oint_C\mathbf B\cdot d\boldsymbol\ell
=\mu_0 I_{\rm wew}+\mu_0\varepsilon_0\frac{d\Phi_E}{dt},\qquad
\nabla\times\mathbf B=\mu_0\mathbf J+\mu_0\varepsilon_0\frac{\partial\mathbf E}{\partial t}.
\]
\(\Phi_E=\int_S\mathbf E\cdot d\mathbf A\), a \(\mathbf J\) to gęstość prądu. Prąd przesunięcia \(I_D=\varepsilon_0\,d\Phi_E/dt\) nie oznacza przepływu ładunków przez dielektryk próżni, lecz wkład zmiennego pola.

Rozważmy idealny kondensator płaski o polu okładek \(A\), odległości \(d\) i pomijalnych efektach brzegowych, ładowany prądem \(I=dQ/dt\). Pomiędzy okładkami \(E=Q/(\varepsilon_0A)\), zatem strumień \(EA=Q/\varepsilon_0\) i \(\varepsilon_0\,d\Phi_E/dt=I\). Powierzchnia rozpięta na tej samej pętli może przecinać przewód z prądem albo przechodzić między okładkami: poprawione prawo daje w obu przypadkach ten sam obieg \(\mathbf B\). Dla kołowej pętli promienia \(r\) między okładkami, przy jednorodnym polu elektrycznym, \(\oint\mathbf B\cdot d\boldsymbol\ell=2\pi rB\), co prowadzi do \(B(r)=\mu_0 I r/(2\pi R^2)\) dla \(r<R\), gdzie \(R\) jest promieniem okładki. Zakładamy ładowanie wolnozmienne i zaniedbujemy pola brzegowe. Poprawka Maxwella jest niezbędna także do przewidzenia fal elektromagnetycznych.

Prąd przesunięcia ma wymiar amperów, ponieważ strumień elektryczny pomnożony przez przenikalność próżni i zróżniczkowany po czasie odpowiada zmianie ładunku w czasie. W przewodniku prąd przewodzenia przenosi ładunki, a między idealnymi okładkami nie ma ich makroskopowego przepływu; mimo to zmieniające się pole tworzy składową w prawie obiegu pola magnetycznego. Założenie jednorodności pola między okładkami jest dobre daleko od ich brzegów, gdzie pole brzegowe można zaniedbać. Jeśli natężenie prądu rośnie lub maleje, zmienia się również tempo zmian strumienia, lecz poprawione prawo obiegu nadal zachowuje ciągłość opisu. W stacjonarnym przypadku \(\partial\mathbf E/\partial t=0\), więc prawo wraca do postaci Ampère’a dla prądu przewodzenia. Odkrycie spójnego wyrazu z pochodną pola nie było tylko poprawką rachunkową: sprzęga ono zmienne pola elektryczne i magnetyczne oraz prowadzi do równań falowych. W ośrodku materialnym używa się często pól \(\mathbf H\) i \(\mathbf D\), a prąd przesunięcia zapisuje się przez zmianę indukcji elektrycznej, z uwzględnieniem odpowiedzi dielektryka. Opis idealnego kondensatora nie obejmuje strat dielektrycznych ani promieniowania.
\subsection{9. Zasada zachowania ładunku z równań Maxwella}
Zakres: lokalne równanie ciągłości dla gęstości ładunku \(\rho\) i gęstości prądu \(\mathbf J\). W próżni różniczkowe prawo Gaussa i prawo Ampère’a--Maxwella mają postać
\[
\nabla\cdot\mathbf E=\rho/\varepsilon_0,\qquad
\nabla\times\mathbf B=\mu_0\mathbf J+\mu_0\varepsilon_0\frac{\partial\mathbf E}{\partial t}.
\]
Bierzemy dywergencję drugiego równania. Dywergencja rotacji dowolnego dostatecznie gładkiego pola jest równa zeru, zatem
\[
0=\mu_0\nabla\cdot\mathbf J+
\mu_0\varepsilon_0\frac{\partial}{\partial t}(\nabla\cdot\mathbf E)
=\mu_0\left(\nabla\cdot\mathbf J+\frac{\partial\rho}{\partial t}\right),
\qquad
\frac{\partial\rho}{\partial t}+\nabla\cdot\mathbf J=0.
\]
To równanie ciągłości: zmiana ładunku w małej objętości jest równa minus strumieniowi prądu wypływającego z niej. Po całkowaniu po ustalonej objętości \(V\) i użyciu twierdzenia Gaussa otrzymujemy \(dQ_V/dt=-\oint_{\partial V}\mathbf J\cdot d\mathbf A\). Ładunek w \(V\) zmienia się więc tylko wtedy, gdy przepływa przez jego brzeg; dla całego izolowanego układu jest zachowany. Gdyby pominąć prąd przesunięcia, dywergencja klasycznego prawa Ampère’a narzucałaby \(\nabla\cdot\mathbf J=0\), co jest sprzeczne z lokalnym gromadzeniem ładunku, np. na okładkach ładowanego kondensatora. Równanie ciągłości jest warunkiem zgodności elektromagnetyzmu z zachowaniem ładunku, a nie dodatkowym założeniem o kształcie trajektorii cząstek. W ośrodkach materialnych podobny bilans obowiązuje przy spójnym rozdzieleniu prądów.

Wyprowadzenie pokazuje, dlaczego poprawka Maxwella jest powiązana z lokalnym bilansem ładunku: pominięcie jej nie pozwalałoby równaniom jednocześnie opisywać gromadzenia się ładunku i spełniać tożsamości wektorowe. Równanie ciągłości ma tę samą strukturę matematyczną co bilans masy w płynach: źródło może wzrosnąć lokalnie tylko wtedy, gdy netto więcej materii napływa niż odpływa. W elektromagnetyzmie ładunek nie powstaje ani nie ginie w zwykłych procesach, lecz może przemieszczać się między fragmentami układu. Kierunek wektora \(d\mathbf A\) jest skierowany na zewnątrz, dlatego dodatni strumień oznacza odpływ prądu i spadek ładunku wewnętrznego. W wersji całkowej objętość może być dowolna, a nie tylko kulista; lokalne równanie wynika z warunku, że bilans ma działać dla każdej małej objętości. W elektrodynamice kwantowej zachowanie ładunku jest również związane z symetrią fazową, lecz tu wystarcza klasyczne wyprowadzenie z równań Maxwella. Przy rozdzielaniu prądów na przewodzenia, polaryzacji i magnetyzacji należy uwzględnić wszystkie składniki zgodnie z przyjętymi polami, aby lokalna zasada ciągłości nadal była spełniona.
\subsection{10. Płaskie fale elektromagnetyczne jako rozwiązania równań Maxwella w próżni}
Zakres: propagacja pól \(\mathbf E,\mathbf B\) w jednorodnej próżni, w obszarze bez ładunków i prądów. Równania Maxwella redukują się tam do \(\nabla\cdot\mathbf E=\nabla\cdot\mathbf B=0\), \(\nabla\times\mathbf E=-\partial_t\mathbf B\) i \(\nabla\times\mathbf B=\mu_0\varepsilon_0\partial_t\mathbf E\). Bierzemy rotację równania Faradaya, stosujemy tożsamość \(\nabla\times(\nabla\times\mathbf E)=\nabla(\nabla\cdot\mathbf E)-\nabla^2\mathbf E\) i podstawiamy prawo Ampère’a--Maxwella. Analogicznie dla \(\mathbf B\):
\[
\nabla^2\mathbf E-\frac{1}{c^2}\frac{\partial^2\mathbf E}{\partial t^2}=0,\qquad
\nabla^2\mathbf B-\frac{1}{c^2}\frac{\partial^2\mathbf B}{\partial t^2}=0,\quad
c=\frac{1}{\sqrt{\mu_0\varepsilon_0}}.
\]
Rozwiązanie płaskie można zapisać \(\mathbf E=\mathbf E_0\cos(\mathbf k\cdot\mathbf r-\omega t)\), z analogicznym \(\mathbf B\), przy \(\omega=c|\mathbf k|\). Równania narzucają \(\mathbf k\cdot\mathbf E_0=\mathbf k\cdot\mathbf B_0=0\), więc oba pola są poprzeczne, oraz
\[
\mathbf B_0=\frac{1}{\omega}\mathbf k\times\mathbf E_0,\qquad
|\mathbf E_0|=c|\mathbf B_0|.
\]
Wektor Poyntinga \(\mathbf S=(\mathbf E\times\mathbf B)/\mu_0\) wskazuje kierunek transportu energii. W próżni nie ma dyspersji w klasycznym modelu: prędkość fazowa jest równa \(c\). Fale niosą energię i pęd oraz mogą przenosić je przez pustą przestrzeń. Przyjęto liniowe równania Maxwella i brak ośrodka; w materii prędkość i relacje amplitudowe zależą od własności materiału.

Zależność dyspersyjna określa, jak częstość fali wiąże się z długością fali; dla próżni \(\omega=c k\), więc wszystkie składowe widma propagują z tą samą prędkością fazową. Kierunek \(\mathbf k\) wskazuje propagację, a wektory \(\mathbf E\) i \(\mathbf B\) leżą w płaszczyźnie poprzecznej do niego i są wzajemnie prostopadłe dla pojedynczej fali płaskiej. Polaryzacja określa kierunek drgań pola elektrycznego, a dla superpozycji składowych można otrzymać polaryzację liniową, kołową lub eliptyczną. Fala harmoniczna o nieskończonym zasięgu jest idealizacją; rzeczywiste źródło wytwarza pakiet o skończonym rozmiarze i czasie trwania. Równania przewidują również odbicie i interferencję, gdy fale napotykają granice, ale ich prosty związek amplitudowy dotyczy obszaru jednorodnego i bezźródłowego. Wektor Poyntinga opisuje chwilowy przepływ energii na jednostkę powierzchni, a jego średnia po okresie daje natężenie fali. Promieniowanie elektromagnetyczne jest więc samopodtrzymującym się zaburzeniem sprzężonych pól, które może rozchodzić się bez ośrodka materialnego. W przybliżeniu geometrzycznym długiej fali kierunek promieni można wiązać z kierunkiem tego przepływu energii.
\subsection{11. Doświadczenia motywujące teorię kwantową: ciało doskonale czarne, fotoefekt, Comptona i widma}
Zakres: porównanie obserwacji z fizyką klasyczną i wskazanie, jaki nowy element wymagał każdy eksperyment. Klasyczny model równowagowego promieniowania ciała doskonale czarnego, oparty na równomiernym podziale energii między mody, prowadził do katastrofy nadfioletowej: nieskończonej energii przy wysokich częstotliwościach. Planck uzyskał widmo zgodne z pomiarem, zakładając wymianę energii w porcjach \(E_n=nh\nu\); energia widma na jednostkę częstotliwości jest proporcjonalna do \(\nu^3/(\exp(h\nu/k_BT)-1)\). W fotoefekcie klasyczna fala nie wyjaśnia progu częstotliwości ani natychmiastowej emisji. Einstein przyjął fotony o energii \(h\nu\), dając równanie \(K_{\max}=h\nu-W\), gdzie \(W\) jest pracą wyjścia; natężenie wpływa głównie na liczbę emitowanych elektronów.

Rozpraszanie Comptona promieniowania rentgenowskiego na elektronach wykazuje zmianę długości fali zależną od kąta: \(\Delta\lambda=(h/m_ec)(1-\cos\vartheta)\). Jest to zgodne z traktowaniem fotonu jako cząstki o pędzie \(h/\lambda\) i zachowaniem energii oraz pędu w zderzeniu, a nie z prostym klasycznym rozpraszaniem fali. Widma atomowe składają się z dyskretnych linii. Model klasyczny przewidywałby ciągłe promieniowanie przyspieszonych ładunków; kwantowo linie odpowiadają różnicom energii stanów: \(h\nu=E_i-E_f\). Każdy przykład pokazuje inny aspekt: kwantowanie wymiany energii, cząstkowe własności światła, pęd fotonu i dyskretność stanów materii. Żaden pojedynczy eksperyment nie jest pełną teorią; ich wspólny opis wymaga mechaniki kwantowej oraz probabilistycznej interpretacji pomiarów.

Szczególnie ważne jest rozróżnienie między energią fotonu a natężeniem wiązki. Zwiększenie częstotliwości zwiększa energię pojedynczego fotonu i może przekroczyć próg emisji, natomiast zwiększenie natężenia przy stałej częstotliwości zwiększa przede wszystkim liczbę fotonów docierających do powierzchni. Zwiększenie temperatury ciała czarnego przesuwa maksimum widma ku wyższym częstotliwościom i zmienia całkowitą moc promieniowania; klasyczna katastrofa nadfioletowa oznaczała błędną rozbieżność, a nie rzeczywiście nieograniczoną emisję. W atomie absorpcja i emisja zachodzą, gdy częstość światła odpowiada różnicy dopuszczonych energii, a nie dowolnej wartości. Zjawiska te nie unieważniają mechaniki klasycznej w jej zakresie stosowalności: klasyczne przybliżenia odzyskuje się zwykle dla dużych liczb kwantowych lub małych kwantów w porównaniu ze skalą energii układu. Eksperymenty stanowiły wskazówki do nowych postulatów, a współczesna teoria wyjaśnia je wspólnie, łącząc dyskretne wymiany energii z falowym opisem prawdopodobieństwa. W każdym przypadku trzeba wskazać, jaka obserwowana wielkość zależy od częstotliwości, a jaka od intensywności lub warunków źródła.
\subsection{12. Zasada nieoznaczoności: pary \(x-p\), kąt--moment pędu, energia--czas i zakres stosowalności}
Zakres: ograniczenia rozrzutów wynikające ze struktury stanów kwantowych, a nie niedoskonałości aparatury. Dla dwóch obserwabli \(A,B\) obowiązuje relacja Robertsona \(\Delta A\,\Delta B\geq|\langle[A,B]\rangle|/2\), gdzie \(\Delta A\) jest odchyleniem standardowym, a \([A,B]=AB-BA\). Ponieważ \([x,p_x]=i\hbar\), mamy \(\Delta x\,\Delta p_x\geq\hbar/2\). Stan skupiony w położeniu musi więc mieć szeroki rozkład pędu; paczka falowa swobodnej cząstki rozpływa się w czasie. To własność stanu, nie zakaz dokładnego pomiaru jednej wielkości.

Dla kąta azymutalnego \(\phi\) i składowej \(L_z\) należy zachować ostrożność: formalna relacja komutacyjna prowadzi do ograniczenia proporcjonalnego do \(\hbar|\langle L_z\rangle|\), lecz kąt jest zmienną okresową i jego wariancja nie zawsze jest standardową miarą rozrzutu. Dokładne ujęcia używają na przykład funkcji sinus i cosinus kąta albo odpowiednio zdefiniowanej wariancji kątowej. Relacja energia--czas ma odmienny status, bo czas w zwykłej mechanice kwantowej jest parametrem, a nie operatorem kanonicznym. Dla stanu o stałym rozrzucie energii skala czasu zmiany obserwabli spełnia ograniczenie Mandelstama--Tammsa \(\tau\Delta E\geq\hbar/2\), przy określonej definicji czasu ewolucji. Nie należy bezwarunkowo traktować jej jak relacji \([x,p]\).

Relacje nieoznaczoności są istotne w mikroświecie, interferencji, wiązaniu stanów i ograniczeniu rozdzielczości fazowej. W makroskopowych układach kwantowych również obowiązują, lecz często są praktycznie niewidoczne wobec dekoherencji i dużych skal działania. Nie zastępują niepewności statystycznej ani błędów przyrządu. Stan minimalnie nieoznaczony, taki jak gaussowska paczka dla \(x,p\), może osiągać równość w standardowej relacji.

Równość w relacji dla położenia i pędu osiągają odpowiednio dobrane stany gaussowskie, ale ogólny stan może mieć iloczyn niepewności znacznie większy od dolnej granicy. Niepewność jest tu rozrzutem wyników w wielu identycznie przygotowanych układach, a nie prostym opisem tego, jak bardzo obserwator ingeruje w pojedynczy pomiar. Dla składowej momentu pędu o niezerowej wartości średniej relacja kątowa wskazuje, że nie można jednocześnie przygotować stanu o ostro określonej orientacji i ostro określonym momencie pędu osiowym. Jednak stan o zerowej wartości średniej może sprawić, że naiwna granica komutatorowa staje się trywialna; dlatego okresowość kąta wymaga starannie zdefiniowanego formalizmu. W relacji energia--czas należy jawnie wskazać, jaką obserwablę i jaką skalę jej zmiany opisuje czas \(\tau\); nie istnieje w standardowym formalizmie uniwersalny operator czasu analogiczny do położenia. W zastosowaniach atomowych nieoznaczoność pędu pomaga wyjaśnić, dlaczego elektron nie może klasycznie spoczywać dokładnie w jądrze bez nieograniczonej energii kinetycznej. Makroskopowy charakter układu nie usuwa zasady, lecz duże wartości działania i oddziaływanie z otoczeniem sprawiają, że rozrzuty względne bywają pomijalne.
\subsection{13. Tunelowanie: stany stacjonarne na prostokątnej barierze \(E<V_0\) i przykład fizyczny}
Zakres: jednowymiarowa cząstka nierelatywistyczna o masie \(m\) w potencjale \(V(x)=V_0\) dla \(0<x<a\) i \(V=0\) poza barierą, przy energii \(0<E<V_0\). Rozwiązujemy równanie stacjonarne Schrödingera
\[
-\frac{\hbar^2}{2m}\frac{d^2\psi}{dx^2}+V(x)\psi=E\psi.
\]
W obszarach swobodnych \(x<0\) i \(x>a\) rozwiązania są kombinacjami fal biegnących o liczbie falowej \(k=\sqrt{2mE}/\hbar\). W barierze równanie przyjmuje postać \(\psi''=\kappa^2\psi\), gdzie \(\kappa=\sqrt{2m(V_0-E)}/\hbar\); rozwiązanie jest kombinacją \(e^{\kappa x}\) i \(e^{-\kappa x}\), a nie propagującą sinusoidą. Ciągłość \(\psi\) i \(\psi'\) na obu granicach pozwala wyznaczyć amplitudy odbicia i transmisji. Dla prostokątnej bariery transmisja wynosi
\[
T=\left[1+\frac{V_0^2\sinh^2(\kappa a)}{4E(V_0-E)}\right]^{-1}.
\]
Choć klasycznie cząstka o tej energii nie przechodzi, kwantowo \(T>0\); dla szerokiej/wysokiej bariery dominuje \(T\sim e^{-2\kappa a}\). Warunki brzegowe wyrażają ciągłość funkcji falowej i jej pochodnej dla skończonego skoku potencjału. Współczynnik odbicia i transmisji sumują się do jedności, gdy nie ma absorpcji. Realny przykład to rozpad alfa: cząstka alfa uwięziona w jądrze tuneluje przez barierę kulombowską, co umożliwia emisję i wyjaśnia silną zależność czasu życia od energii cząstki. Model jednej prostokątnej bariery jest ilustracją mechanizmu; w jądrze rzeczywisty potencjał jest przestrzennie zmienny i oblicza się transmisję metodami przybliżonymi, np. WKB.

Fizycznie współczynnik transmisji jest stosunkiem strumienia prawdopodobieństwa po stronie wyjściowej do strumienia fali padającej, a nie po prostu kwadratem dowolnego stosunku amplitud. W tym przypadku fale wejściowa i wyjściowa mają tę samą energię i liczbę falową, co upraszcza taki stosunek. W obliczeniu trzeba uwzględnić falę odbitą przed barierą, ponieważ dopasowanie na dwóch granicach tworzy interferencję wewnątrz bariery. Szerokość bariery wpływa wykładniczo na prawdopodobieństwo przejścia, zatem niewielka zmiana szerokości może radykalnie zmienić częstość procesu. Ten mechanizm wykorzystuje mikroskop skaningowy tunelowy: elektrony tunelują między ostrzem a powierzchnią, a prąd zależy bardzo silnie od odległości, dzięki czemu można obrazować strukturę powierzchni. W przykładzie alfa energia całkowita jest mniejsza od maksimum bariery kulombowskiej, lecz funkcja falowa nie znika nagle na jej granicy. Warunki ciągłości zapewniają poprawne połączenie rozwiązań w każdym regionie; delta-potencjał lub nieskończona ściana wymagałyby osobnego omówienia warunków brzegowych. Tunelowanie nie oznacza klasycznego przejścia cząstki przez barierę po trajektorii z chwilowo ujemną energią kinetyczną; opisuje niezerową amplitudę transmisji kwantowej.
\subsection{14. Równanie Schrödingera i zachowanie cząstki swobodnej}
Zakres: nierelatywistyczny opis kwantowy cząstki masy \(m\), bez spinu w podstawowym modelu, poruszającej się w potencjale \(V(\mathbf r,t)\). Równanie zależne od czasu wynika z postulatu ewolucji amplitudy \(\psi\):
\[
i\hbar\frac{\partial\psi}{\partial t}
=\left[-\frac{\hbar^2}{2m}\nabla^2+V(\mathbf r,t)\right]\psi.
\]
Operator \(-i\hbar\nabla\) odpowiada pędowi, a \(-\hbar^2\nabla^2/(2m)\) energii kinetycznej. Gęstość prawdopodobieństwa wynosi \(|\psi|^2\); przy normalizacji \(\int|\psi|^2d^3r=1\) całkowite prawdopodobieństwo jest jednością. Dla rzeczywistego potencjału równanie prowadzi do równania ciągłości prawdopodobieństwa z prądem \(\mathbf j=(\hbar/m)\operatorname{Im}(\psi^*\nabla\psi)\), co odzwierciedla zachowanie normy.

Dla cząstki swobodnej \(V=0\) rozwiązaniami są fale płaskie
\[
\psi(\mathbf r,t)=A e^{i(\mathbf k\cdot\mathbf r-\omega t)},\quad
\mathbf p=\hbar\mathbf k,\quad E=\hbar\omega=\frac{\hbar^2k^2}{2m}.
\]
Fala płaska ma dokładnie określony pęd, lecz nie jest normalizowalna w całej przestrzeni; fizyczne stany tworzy się jako paczki falowe, czyli superpozycje fal o różnych \(\mathbf k\). Dla pakietu o rozrzucie pędu położenie może być początkowo skupione, a następnie dyspersja relacji \(E=p^2/(2m)\) powoduje jego rozprzestrzenianie. Równanie stacjonarne \(H\phi=E\phi\) otrzymujemy przy potencjale niezależnym od czasu i rozdzieleniu \(\psi=\phi e^{-iEt/\hbar}\). Teoria zakłada prędkości dużo mniejsze od \(c\), a przy niskich energiach nie opisuje kreacji cząstek ani efektów relatywistycznych.

Interpretacja probabilistyczna wymaga odróżnienia samej funkcji falowej od wielkości obserwowalnych: globalny czynnik fazowy nie zmienia gęstości prawdopodobieństwa, natomiast faza względna między składowymi superpozycji decyduje o interferencji. W równaniu Schrödingera Hamiltonian jest operatorem energii, a jego postać zależy od potencjału i warunków fizycznych. Dla swobodnej cząstki średnia prędkość paczki jest związana z pędem, lecz faza płaskiej fali rozchodzi się z prędkością fazową różną od prędkości grupowej; informacja o ruchu pakietu jest związana z prędkością grupową. Jeżeli cząstka znajduje się w skończonym obszarze, dozwolone stany i energie zależą od warunków brzegowych nawet wtedy, gdy potencjał wewnątrz jest stały. Równanie stacjonarne jest wygodne dla poziomów energii, lecz samo nie opisuje superpozycji o różnych energiach jako jednego stanu o stałej fazie czasowej. W doświadczeniu pomiar położenia daje pojedynczy wynik, a rozkład wielu powtórzeń określa \(|\psi|^2\). Równanie jest liniowe, więc suma rozwiązań także jest rozwiązaniem; to matematyczna podstawa interferencji kwantowej. Gdy pole zewnętrzne zależy od czasu, energia układu nie musi być zachowana, choć ewolucja nadal zachowuje normę.
\subsection{15. Funkcje falowe układów wielu cząstek: bozony, fermiony i przykłady}
Zakres: stan kwantowy \(N\) cząstek opisany funkcją \(\Psi(\mathbf r_1,\ldots,\mathbf r_N)\), z dodatkowymi zmiennymi spinu, jeśli są potrzebne. Dla cząstek identycznych zamiana etykiet nie może prowadzić do nowego fizycznie rozróżnialnego stanu. Funkcja całkowita (przestrzeń i spin) jest symetryczna dla bozonów i antysymetryczna dla fermionów:
\[
\Psi(\ldots,\mathbf r_i,\ldots,\mathbf r_j,\ldots)
=\begin{cases}
+\Psi(\ldots,\mathbf r_j,\ldots,\mathbf r_i,\ldots)&\text{bozony},\\
-\Psi(\ldots,\mathbf r_j,\ldots,\mathbf r_i,\ldots)&\text{fermiony}.
\end{cases}
\]
Antysymetria daje zakaz Pauliego: dwa identyczne fermiony nie mogą zajmować tego samego stanu jednocząstkowego, bo wymiana dwóch identycznych kolumn wyznacznika Slatera zmienia znak, a przy powtórzeniu stanu wyznacznik znika. Bozony mogą zajmować wspólny stan, co umożliwia kondensację Bosego--Einsteina i koherentne stany światła. Elektrony, protony i neutrony są fermionami; fotony oraz atomy helu-4 zachowują się jak bozony w odpowiednim zakresie energii.

Dla niezależnych cząstek funkcje jednocząstkowe można łączyć iloczynem dla rozróżnialnych obiektów, symetryzować lub antysymetryzować dla identycznych. Dla oddziałujących cząstek stan zwykle nie rozkłada się na prosty iloczyn, co prowadzi do splątania. Warto rozdzielać spin i część przestrzenną: ich łączna symetria musi odpowiadać statystyce cząstek. Na przykład dwa elektrony w stanie singletowym spinowym mają antysymetryczną część spinową, więc część przestrzenna może być symetryczna. Prawdopodobieństwa wyników pomiaru wynikają z \(|\Psi|^2\), ale amplitudy dla cząstek identycznych interferują po zamianie etykiet. W makroskopowych gazach rozróżnienie statystyk staje się ważne, gdy długość fali de Broglie’a cząstek jest porównywalna z odległością między nimi.

W układzie wielu identycznych cząstek etykiety \(i\) i \(j\) są pomocnicze, dlatego fizyczny stan nie może zależeć od tego, którą cząstkę nazwano pierwszą. Symetryzacja amplitud prowadzi do efektów wymiany nawet wtedy, gdy cząstki nie oddziałują siłą; dla elektronów daje to między innymi strukturę powłok atomowych. Zasada Pauliego dotyczy pełnego stanu jednocząstkowego, obejmującego także spin, dlatego dwa elektrony mogą mieć tę samą funkcję przestrzenną, jeśli zajmują różne stany spinowe. W przewodniku wiele fermionów zapełnia stany do energii Fermiego, natomiast fotony w rezonatorze mogą gromadzić się w tym samym modzie. Atomy helu-4 są bozonami złożonymi z parzystej liczby składników fermionowych; ich zachowanie bozonowe ma sens, gdy rozdzielczość doświadczenia nie ujawnia struktury wewnętrznej. Dla dwóch cząstek należy stosować odpowiednią kombinację symetryczną lub antysymetryczną, a nie tylko pojedynczy iloczyn funkcji. W konsekwencji rozkłady prawdopodobieństwa mogą mieć wzmocnienie lub osłabienie przy małych odległościach, znane odpowiednio jako grupowanie bozonowe i antygrupowanie fermionowe. Statystyka wpływa na równowagowe własności materii, choć dla rozrzedzonego i gorącego gazu klasyczne przybliżenie cząstek rozróżnialnych bywa wystarczające.
\subsection{16. Opis mikrostanowy termodynamiki oraz pierwsza i druga zasada}
Zakres: połączenie makroskopowych zmiennych stanu z mikroskopowymi konfiguracjami cząstek. Mikrostan określa pełną konfigurację układu, klasycznie przez położenia i pędy wszystkich cząstek, a kwantowo przez stan układu. Makrostan określają wielkości takie jak energia, objętość, liczba cząstek, temperatura i ciśnienie; wiele mikrostanów może odpowiadać temu samemu makrostanowi. W równowadze makrostan jest związany z liczbą dostępnych mikrostanów \(\Omega\), a entropia Boltzmanna to \(S=k_B\ln\Omega\) dla równoprawdopodobnych mikrostanów. W ujęciu probabilistycznym \(S=-k_B\sum_i p_i\ln p_i\).

Pierwsza zasada jest bilansem energii. Przy konwencji, że \(W\) oznacza pracę wykonaną przez układ, \(\Delta U=Q-W\), gdzie \(U\) jest energią wewnętrzną, a \(Q\) ciepłem dostarczonym układowi. Dla pracy objętościowej \(W=\int p_{\rm ext}\,dV\); w procesie odwracalnym quasi-statycznym \(p_{\rm ext}=p\). Druga zasada ogranicza kierunek procesów: dla układu izolowanego \(\Delta S\geq0\), a równość zachodzi w granicy procesu odwracalnego. Mikroskopowe równania ruchu mogą być odwracalne, lecz typowy makroproces spontanicznie prowadzi do makrostanów odpowiadających znacznie większej liczbie mikrostanów. Przykładem jest swobodne rozprężanie gazu, w którym gaz wypełnia większą objętość, choć każda trajektoria mikroskopowa jest zgodna z mechaniką. Dla wymiany ciepła z otoczeniem \(\Delta S_{\rm całk}\geq0\); entropia samego układu może maleć, jeśli otoczenie zwiększy ją bardziej. Termodynamika równowagowa zakłada makroskopowy opis i nie opisuje szczegółowo czasu relaksacji.

W rachunku termodynamicznym trzeba wyraźnie ustalić znaki: jeśli do układu doprowadzamy ciepło, \(Q\) jest dodatnie, a jeśli układ wykonuje pracę nad otoczeniem, \(W\) jest dodatnie. Dzięki temu pierwsza zasada zachowuje tę samą treść przy rozprężaniu, sprężaniu i ogrzewaniu. Entropia nie oznacza potocznego nieporządku, lecz ilościowe zliczenie lub miarę niepewności mikrostanów zgodnych z opisem makroskopowym. Układ o większej objętości ma zwykle więcej dostępnych położeń cząstek, dlatego wzrasta liczba mikrostanów i entropia gazu podczas swobodnego rozprężania. Odwrócenie takiego procesu nie jest zakazane przez bilans energii, ale wymaga szczególnego przygotowania mikroskopowego i nie zachodzi spontanicznie w typowym makroskopowym układzie. W procesie odwracalnym można zdefiniować zmianę entropii przez wymianę ciepła podzieloną przez temperaturę, ale dla procesu nieodwracalnego tę relację stosuje się do drogi odwracalnej łączącej te same stany, nie do rzeczywistego strumienia. Pierwsza zasada jest bilansem ilości energii, druga zaś określa kierunek przemian i ogranicza sprawność silników cieplnych. Obie zasady są komplementarne i żadna nie wynika z samej drugiej.
\subsection{17. Zespół kanoniczny klasycznych cząstek: funkcja podziału i wielkości makroskopowe}
Zakres: układ w równowadze termicznej z dużym rezerwuarem o temperaturze \(T\), o ustalonej liczbie cząstek \(N\) i objętości \(V\). Mikrostan o energii \(E_i\) ma prawdopodobieństwo \(p_i=e^{-\beta E_i}/Z\), gdzie \(\beta=1/(k_BT)\), a suma statystyczna
\[
Z=\sum_i e^{-\beta E_i}
\]
normalizuje rozkład. Dla klasycznego gazu Hamiltonianu \(H(\mathbf p,\mathbf q)\) suma zastępuje się całką fazową. Dla \(N\) identycznych, nierelatywistycznych, nieoddziałujących cząstek w pojemniku:
\[
Z_N=\frac{1}{N!h^{3N}}\int e^{-\beta H}\,d^{3N}q\,d^{3N}p
=\frac{1}{N!}\left(\frac{V}{\lambda_T^3}\right)^N,\quad
\lambda_T=\frac{h}{\sqrt{2\pi m k_BT}}.
\]
Czynnik \(1/N!\) uwzględnia nierozróżnialność klasycznych cząstek, a \(\lambda_T\) jest termiczną długością fali de Broglie’a. Przybliżenie klasyczne wymaga małego parametru gęstości \(n\lambda_T^3\ll1\).

Energia średnia i energia swobodna Helmholtza wynikają z \(Z\):
\[
U=-\frac{\partial\ln Z}{\partial\beta},\qquad
F=-k_BT\ln Z,\qquad
p=-\left(\frac{\partial F}{\partial V}\right)_{T,N}.
\]
Dla gazu doskonałego daje to równanie stanu \(pV=Nk_BT\) oraz \(U=3Nk_BT/2\). Entropię wyznacza \(S=-\left(\partial F/\partial T\right)_{V,N}\), a fluktuacja energii spełnia \(\langle(\Delta E)^2\rangle=k_BT^2C_V\), gdzie \(C_V=(\partial U/\partial T)_V\). Tak więc jedna funkcja \(Z\) koduje średnie wielkości makroskopowe i fluktuacje. W obecności oddziaływań trzeba użyć pełnego \(H\), a obliczenie całki może wymagać przybliżeń. Zespół kanoniczny opisuje wymianę energii z termostatem; energia nie jest ustalona w każdym mikroświecie, choć jej średnia jest określona.
Wzór na funkcję podziału rozdziela całkę pędową od położeniowej, ponieważ Hamiltonian gazu doskonałego jest sumą energii kinetycznych, a potencjał wewnątrz pojemnika jest stały. Całka Gaussowska po pędach daje czynnik związany z długością termiczną, zaś całka po położeniach daje potęgę objętości. Czynnik \(N!\) usuwa sztuczne liczenie różnych permutacji identycznych cząstek jako odrębnych mikrostanów i jest istotny dla ekstensywności entropii. Pochodna \(\ln Z\) po \(\beta\) waży każdy stan jego energią, dlatego daje średnią energię, a pochodne wyższych rzędów opisują jej fluktuacje. Zależność od objętości prowadzi do ciśnienia jako siły uogólnionej sprzężonej z objętością; w gazie doskonałym ta metoda odtwarza prawo empiryczne. Przy zwiększaniu liczby cząstek fluktuacje energii rosną, ale ich względny udział zwykle maleje, więc wielkości makroskopowe stają się stabilne. W pobliżu przejść fazowych fluktuacje mogą jednak być duże, a proste przybliżenie gazu nie wystarcza. Zespół kanoniczny jest naturalny dla próbki utrzymywanej w kontakcie cieplnym z otoczeniem, natomiast układ izolowany o ściśle ustalonej energii opisuje zespół mikrokanoniczny.

\clearpage
\section{Ekonofizyka}
\subsection{1. Konstrukcja i analiza portfela papierów wartościowych}

Portfel jest zbiorem pozycji w instrumentach finansowych, opisanym przez udziały \(w_i\), gdzie \(w_i\) oznacza część kapitału przypisaną do \(i\)-tego składnika. Przy pełnym zainwestowaniu i braku krótkiej sprzedaży zachodzi \(\sum_{i=1}^{N}w_i=1\) oraz \(w_i\geq 0\). Niech \(P_{i,t}\) oznacza cenę instrumentu w chwili \(t\), a \(R_{i,t+1}=(P_{i,t+1}-P_{i,t})/P_{i,t}\) jego prostą stopę zwrotu. Dla składu ustalonego na początku okresu stopa zwrotu portfela wynosi \(R_{p,t+1}=\sum_i w_iR_{i,t+1}\). Dla stóp logarytmicznych \(r_{i,t+1}=\ln(P_{i,t+1}/P_{i,t})\) ta liniowa zależność jest przybliżeniem, dokładnym w granicy małych zmian cen.

Analiza zaczyna się od oszacowania średnich \(\mu_i=\mathbb E[R_i]\) i macierzy kowariancji \(\Sigma\), której element \(\Sigma_{ij}=\operatorname{Cov}(R_i,R_j)\) opisuje współzmienność pary składników. Oczekiwana stopa zwrotu oraz wariancja portfela to odpowiednio
\[
\mathbb E[R_p]=\mathbf w^{\mathsf T}\boldsymbol\mu,\qquad
\operatorname{Var}(R_p)=\mathbf w^{\mathsf T}\Sigma\mathbf w.
\]
Przykładowo, dwa aktywa o jednakowych wariancjach mogą tworzyć portfel o mniejszej wariancji niż każde z osobna, jeśli ich korelacja jest mniejsza od jedności. W praktyce szeregi cen należy uzgodnić co do częstotliwości, waluty, dywidend i brakujących obserwacji, a stopy zwrotu obliczać bez wprowadzania informacji z przyszłości. Należy też rozróżniać wynik historyczny od oczekiwanego wyniku: średnie i kowariancje z próby są obarczone błędem estymacji, w szczególności przy krótkich szeregach lub wielu aktywach. Opis ten jest narzędziem matematycznym do analizy ryzyka i nie stanowi rekomendacji inwestycyjnej.

W konkretnym przykładzie portfel złożony w równych częściach z dwóch aktywów ma zwrot równy średniej ich zwrotów, gdy udziały pozostają stałe w danym okresie. Jeśli ceny składników zmieniają się różnie, ich udziały rynkowe odchylają się od początkowych i rebalansowanie przywraca wybrany skład; transakcje te mogą jednak generować koszty oraz podatki. Trzeba jawnie określić, czy model zakłada krótką sprzedaż, możliwość lokowania gotówki i stałe wagi, ponieważ każda z tych decyzji zmienia dopuszczalny zbiór portfeli. Miary ryzyka zależą ponadto od horyzontu pomiaru, a przeskalowanie zmienności pierwiastkiem z czasu wymaga dodatkowych założeń o zależności zwrotów. Dlatego analiza portfela powinna rozpoczynać się od precyzyjnego opisania celu i ograniczeń.

\subsection{2. Portfel minimalnej wariancji i portfel optymalny}

Portfel minimalnej wariancji to taki dopuszczalny wektor wag \(\mathbf w\), dla którego \(\mathbf w^{\mathsf T}\Sigma\mathbf w\) jest najmniejsze przy zadanych ograniczeniach. Przy dopuszczeniu krótkiej sprzedaży, warunku \(\mathbf 1^{\mathsf T}\mathbf w=1\) i dodatnio określonej macierzy \(\Sigma\), bez dodatkowych ograniczeń rozwiązanie ma postać
\[
\mathbf w_{\mathrm{MV}}=
\frac{\Sigma^{-1}\mathbf 1}{\mathbf 1^{\mathsf T}\Sigma^{-1}\mathbf 1}.
\]
Wynika ono z minimalizacji formy kwadratowej przy ograniczeniu sumy wag: lagranżjan \(\mathcal L=\mathbf w^{\mathsf T}\Sigma\mathbf w-\lambda(\mathbf 1^{\mathsf T}\mathbf w-1)\) daje \(2\Sigma\mathbf w=\lambda\mathbf 1\); podstawienie warunku sumy wyznacza stałą i powyższy wynik. Przy zakazie krótkiej sprzedaży albo limitach wag rozwiązanie może leżeć na brzegu obszaru dopuszczalnego i trzeba rozwiązać zadanie z nierównościami.

Portfel optymalny zależy od kryterium: może maksymalizować oczekiwany zwrot przy ograniczeniu ryzyka, minimalizować ryzyko przy wymaganym zwrocie albo maksymalizować użyteczność. W klasycznym ujęciu średnia–wariancja dla celu \(\boldsymbol\mu^{\mathsf T}\mathbf w=\mu_0\) rozwiązuje się minimalizując wariancję pod warunkami \(\mathbf 1^{\mathsf T}\mathbf w=1\) i tym warunkiem zwrotu. Zbiór takich rozwiązań tworzy granicę efektywną: dla każdego poziomu wariancji nie istnieje portfel o większym średnim zwrocie. „Optymalność” nie jest zatem własnością niezależną od założeń i preferencji. Estymacja \(\boldsymbol\mu\) bywa szczególnie niestabilna, a odwracanie źle uwarunkowanej \(\Sigma\) może prowadzić do skrajnych wag. Wyniki należy interpretować warunkowo względem modelu, okresu próby, kosztów transakcyjnych i ograniczeń; modelowa granica nie zapewnia osiągnięcia przewidywanych wyników.

Prosty przykład pokazuje rolę korelacji: dla dwóch aktywów o odchyleniu standardowym \(\sigma\), jednakowych oczekiwanych zwrotach i korelacji \(\rho=0{,}2\), portfel z równymi wagami ma wariancję \(\sigma^2(1+\rho)/2=0{,}6\sigma^2\). Jest to mniej niż wariancja każdego aktywa, równa \(\sigma^2\), ale tylko przy przyjętych parametrach i w rozpatrywanym okresie. Jeśli zwroty oczekiwane są różne, portfel minimalnej wariancji nie musi spełniać wymaganego celu zwrotu, więc należy odróżnić go od konkretnego portfela efektywnego. W obliczeniach empirycznych warto porównać wagi przed i po niewielkiej zmianie próby oraz uwzględnić koszty rebalansowania. Takie testy wrażliwości ujawniają, czy pozorna przewaga wynika z założeń, które trudno wiarygodnie oszacować.

\subsection{3. Korelacje cen i stóp zwrotu; taksonomia portfeli}

Korelacja liniowa stóp zwrotu dwóch aktywów definiowana jest jako
\[
\rho_{ij}=\frac{\operatorname{Cov}(R_i,R_j)}
{\sqrt{\operatorname{Var}(R_i)\operatorname{Var}(R_j)}}\in[-1,1].
\]
Wartość bliska \(1\) oznacza silną dodatnią współzmienność liniową, bliska \(-1\) — ujemną, a bliska zeru — brak wyraźnej zależności liniowej; zero nie wyklucza zależności nieliniowej. Dla dwóch aktywów wariancja portfela wynosi \(w^2\sigma_1^2+(1-w)^2\sigma_2^2+2w(1-w)\rho_{12}\sigma_1\sigma_2\). Wyraz z korelacją wyjaśnia dywersyfikację: przy niższej korelacji łączna zmienność może być mniejsza, lecz nie wynika z tego gwarancja zysku ani ochrony przed każdym spadkiem.

Korelacje cenowych poziomów często są mylące: niestacjonarne ceny mogą wykazywać wysoką korelację jedynie z powodu wspólnego trendu. Z tego względu zwykle bada się stopy zwrotu, sprawdzając stacjonarność, autokorelacje i zależności w różnych horyzontach. Korelacja Pearsona jest wrażliwa na obserwacje skrajne i mierzy zależność liniową. Współczynniki rangowe opisują zależność monotoniczną, natomiast zależność ogonową można analizować metodami kopulowymi lub prawdopodobieństwami wspólnych ekstremów. Korelacje mogą zmieniać się w czasie, zwłaszcza w okresach stresu rynkowego; pojedyncza macierz z całej próby może to ukryć.

W taksonomii portfeli można rozróżnić portfele rynkowe, branżowe, geograficzne, walutowe, obligacyjne, akcyjne i mieszane — według rodzaju lub ekspozycji składników. Według metody konstrukcji wyróżnia się portfel ważony kapitalizacją, równoważony, minimalnej wariancji i oparty na czynnikach ryzyka. Są to różne klasyfikacje: portfel branżowy może być równoważony lub ważony kapitalizacją. Należy podać kryterium klasyfikacji, okres i definicję zwrotu, aby porównanie miało sens; etykieta nie określa sama w sobie zachowania ryzyka.

Przy estymacji korelacji wynik zależy od długości i położenia okna obserwacji. Okno krótkie szybciej odzwierciedla zmiany, lecz zwiększa szum estymatora; okno długie daje stabilniejsze oszacowanie, ale może uśredniać kilka odmiennych reżimów. Macierz korelacji powinna być symetryczna i dodatnio półokreślona, co jest ważne przy wykorzystaniu jej w dalszych obliczeniach portfelowych. Braki danych i różne godziny zamknięcia rynków mogą sztucznie zaniżać lub zniekształcać współzależności. Przy porównywaniu klasyfikacji portfeli należy więc stosować porównywalne próbki i uwzględniać niepewność oszacowań. Zmiana korelacji w kryzysie jest możliwą własnością danych, a nie prawem, które musi powtórzyć się w każdym epizodzie.

\subsection{4. Krach finansowy jako przejście fazowe: opis, prognozowanie i metody badania}

W analogii do przejścia fazowego krach opisuje się jako gwałtowne, skoordynowane przejście układu rynkowego od stanu względnie spokojnego do stanu dużych spadków i podwyższonej zmienności. Elementami układu mogą być inwestorzy, a oddziaływania społeczne, naśladownictwo i ograniczenia płynności sprzyjają korelacjom między decyzjami. W języku fizyki „parametrem porządku” bywa zagregowany kierunek zmian lub poziom współzależności, a „parametrem kontrolnym” — zadłużenie, płynność, niepewność bądź dźwignia. Są to analogie modelowe, nie dowód, że rynek jest układem termodynamicznym ani że istnieje uniwersalny punkt krytyczny.

Ilościowy opis krachu może obejmować logarytmiczną cenę \(p(t)=\ln P(t)\), stopę zwrotu \(r_t=p(t)-p(t-\Delta t)\), zmienność zrealizowaną, obsunięcie od lokalnego maksimum, miary płynności i strukturę korelacji. Jedna hipoteza przedkrytyczna zakłada przyspieszające oscylacje log-okresowe:
\[
p(t)=A+B(t_c-t)^m+C(t_c-t)^m
\cos\!\bigl(\omega\ln(t_c-t)+\phi\bigr),\quad t<t_c,
\]
gdzie \(t_c\) jest hipotetycznym czasem osobliwości, a parametry dopasowuje się do danych. Model może opisywać pewne epizody, lecz dopasowanie po fakcie nie daje wiarygodnej prognozy konkretnej daty krachu.

Badania obejmują analizę zdarzeń, testy zmian reżimu, symulacje agentowe, modele sieciowe oraz modele zmienności i ekstremów. Prognozowanie powinno być probabilistyczne, z walidacją poza próbą, porównaniem z prostymi punktami odniesienia i oceną fałszywych alarmów. Krache są rzadkie, ich definicje zależą od progu i horyzontu, a parametry zmieniają się w czasie; ponadto reżimy mogą zmienić się nagle bez czytelnych sygnałów wyprzedzających. Zależności i sygnały, które wydają się stabilne w danych historycznych, nie muszą przetrwać nowego kryzysu. Analogie fizyczne są użyteczne do formułowania hipotez, ale nie powinny być przedstawiane jako pewna metoda przewidywania ani jako podstawa gwarantowanych decyzji.

\subsection{5. Statystyczne własności rozkładów stóp zwrotu}

Rozkład stóp zwrotu opisuje prawdopodobieństwo wielkości zmian cen w określonym horyzoncie. W praktyce często bada się \(r_t=\ln(P_t/P_{t-1})\), a nie poziomy \(P_t\), ponieważ zwroty są zwykle bliższe stacjonarności. Średnia i wariancja to podstawowe parametry, ale nie wystarczają do pełnego opisu. Standaryzowana skośność \(S=\mathbb E[(r-\mu)^3]/\sigma^3\) mierzy asymetrię, a kurtoza \(\kappa=\mathbb E[(r-\mu)^4]/\sigma^4\) określa względną wagę ogonów; kurtoza nadmiarowa to \(\kappa-3\), równa zeru dla rozkładu normalnego. W finansowych próbach często obserwuje się skupianie dużych wartości bezwzględnych zwrotów i ogony grubsze niż normalne, a także zależność tych cech od rynku i skali czasu.

Przy założeniu \(r_t\sim\mathcal N(\mu,\sigma^2)\), prawdopodobieństwo zwrotu mniejszego od progu \(a\) można obliczyć jako \(\Phi((a-\mu)/\sigma)\). Jeśli rzeczywisty rozkład ma cięższe ogony, takie obliczenie może zaniżać prawdopodobieństwo skrajnych zmian. Należy więc badać histogramy i wykresy kwantyl–kwantyl, estymować empiryczne kwantyle, testować normalność oraz uwzględniać niepewność próby. Sama skośność lub kurtoza jest niestabilna przy krótkich szeregach, bo momenty wysokiego rzędu silnie reagują na pojedyncze ekstremum.

Rozkład bezwarunkowy łączy obserwacje pochodzące z różnych poziomów zmienności. Nawet przy warunkowo normalnych innowacjach zmienność zmieniająca się w czasie może tworzyć bezwarunkowo grube ogony. Dlatego warto rozdzielać rozkład zwrotu od rozkładu innowacji po standaryzacji przez oszacowaną zmienność. Wyniki zależą też od częstotliwości, sposobu agregacji, mikrostruktury rynku i zdarzeń wyjątkowych; dzienne zwroty nie muszą mieć tych samych własności co minutowe lub miesięczne. Modele ogonów, takie jak rozkłady \(t\)-Studenta lub uogólnione rozkłady Pareto dla nadwyżek ponad próg, wymagają sprawdzenia dopasowania i analizy wrażliwości progu. Rozkład historyczny nie wyznacza dokładnie rozkładu przyszłego.

\subsection{6. Modelowanie zwrotów przez AR, MA i ARMA a rozkład Gaussa}

Modele AR, MA i ARMA opisują zależność liniową w szeregu stacjonarnym \(x_t\), którym może być stopa zwrotu po odjęciu średniej. Model AR(\(p\)) ma postać \(x_t=c+\sum_{i=1}^{p}\phi_i x_{t-i}+\varepsilon_t\), gdzie \(\varepsilon_t\) jest białym szumem. W modelu MA(\(q\)) \(x_t=\mu+\varepsilon_t+\sum_{j=1}^{q}\theta_j\varepsilon_{t-j}\); bieżąca obserwacja zależy od bieżącego i opóźnionych szoków. Model ARMA(\(p,q\)) łączy oba mechanizmy. Stacjonarność AR wymaga, aby pierwiastki wielomianu autoregresji leżały poza okręgiem jednostkowym; odwracalność MA analogicznie dotyczy wielomianu MA. Warunki te zapewniają stabilną interpretację parametrów i reprezentację procesu.

Rząd modelu można wstępnie oceniać przez funkcje autokorelacji ACF i PACF, a następnie dobierać kryterium informacyjne, na przykład AIC, oraz badać reszty. Po oszacowaniu modelu reszty powinny nie wykazywać pozostałej liniowej autokorelacji; diagnostyka obejmuje także zmienność reszt i obserwacje odstające. W przypadku stacjonarnego AR(1), \(x_t=\phi x_{t-1}+\varepsilon_t\), z \(|\phi|<1\), niezależnymi \(\varepsilon_t\sim\mathcal N(0,\sigma_\varepsilon^2)\), proces jest Gaussowski: \(x_t\) jest sumą ważoną niezależnych zmiennych normalnych. Jego wariancja wynosi \(\sigma_\varepsilon^2/(1-\phi^2)\), a autokorelacja rzędu \(k\) — \(\phi^k\). Ogólniej liniowy model ARMA z gaussowskimi innowacjami ma wspólnie gaussowski rozkład.

Sama etykieta ARMA nie oznacza normalności: przy nienormalnych innowacjach rozkład może być nienormalny, nawet jeśli zależności liniowe są prawidłowo opisane. Ponadto zwykły ARMA modeluje warunkową średnią, a nie skupianie zmienności. Zwroty mogą mieć słabą autokorelację, ale silną autokorelację kwadratów; wtedy potrzebny jest model zmienności, np. ARCH/GARCH, lub inna rozszerzona specyfikacja. Dobry wynik dopasowania do próby nie gwarantuje prognoz poza próbą, a testy powinny uwzględniać zmianę reżimów i niepewność doboru modelu.

Warto odróżnić modelowanie średniej od założenia o rozkładzie: równanie ARMA opisuje sposób łączenia bieżących i przeszłych składników, natomiast rozkład innowacji stanowi osobny wybór. Przy niezależnych, gaussowskich innowacjach i stabilnych parametrach liniowa kombinacja daje rozkład normalny, ale w skończonym szeregu wnioski o normalności pozostają empiryczne. Prognozę punktową można uzyskać przez warunkową wartość oczekiwaną, a przedział prognozy wymaga również określenia rozkładu błędu i jego wariancji. Diagnostyka powinna być przeprowadzona na resztach z kolejnych okresów, nie tylko na tych samych danych, które posłużyły do dopasowania. Jeżeli średnia jest bliska zeru i nie ma autokorelacji, ARMA może nie poprawiać prostego modelu stałej średniej; prostszy model warto zachować jako punkt odniesienia.

\subsection{7. Stochastyczne procesy ARCH/GARCH dla finansowych szeregów czasowych}

Modele ARCH/GARCH opisują zmienność warunkową, która zmienia się w czasie i zależy od wcześniejszych szoków. Niech \(r_t=\mu_t+\varepsilon_t\), gdzie \(\varepsilon_t=\sigma_t z_t\), \(z_t\) ma średnią zero, wariancję jeden i jest niezależne w czasie. W ARCH(\(q\)) wariancja warunkowa spełnia
\[
\sigma_t^2=\omega+\sum_{i=1}^{q}\alpha_i\varepsilon_{t-i}^2,
\]
gdzie zwykle \(\omega>0\), \(\alpha_i\geq0\). Duży szok zwiększa prognozowaną wariancję na kolejne okresy, co modeluje skupianie zmienności. GARCH(\(p,q\)) dodaje opóźnione wariancje:
\[
\sigma_t^2=\omega+\sum_{i=1}^{q}\alpha_i\varepsilon_{t-i}^2+
\sum_{j=1}^{p}\beta_j\sigma_{t-j}^2.
\]
Dla GARCH(1,1) przy dodatnich parametrach warunek \(\alpha_1+\beta_1<1\) zapewnia skończoną bezwarunkową wariancję \(\omega/(1-\alpha_1-\beta_1)\), przy standardowych założeniach. Wysoka wartość \(\beta_1\) często oznacza powolny zanik wpływu szoku, choć interpretacja zależy od parametrów i rozkładu innowacji.

Estymacja zwykle polega na maksymalizacji wiarygodności, przyjmując np. normalne albo \(t\)-Studenta \(z_t\). Rozkład \(t\) może lepiej ujmować ciężkie ogony reszt, ale nie naprawi błędnej dynamiki. Po dopasowaniu trzeba analizować standaryzowane innowacje \(\hat z_t=\hat\varepsilon_t/\hat\sigma_t\), ich kwadraty, autokorelacje i dopasowanie ogonów. Jeśli pozostaje skupianie zmienności, specyfikacja jest niewystarczająca. Asymetryczne warianty, takie jak GJR-GARCH, dopuszczają odmienną reakcję na szoki dodatnie i ujemne.

ARCH/GARCH nie gwarantuje trafnego prognozowania krachów: model zakłada określoną dynamikę wariancji i może niedoszacować strukturalnej zmiany reżimu lub wyjątkowej obserwacji. Wyniki zależą od wyboru rozkładu, długości próby i częstotliwości danych. Należy oddzielić prognozę wariancji warunkowej od przewidywania kierunku zwrotu; typowy GARCH modeluje przede wszystkim pierwszą z tych wielkości. Jest to opis statystyczny, a nie ocena bezpieczeństwa konkretnego instrumentu.

Kluczowe rozróżnienie dotyczy wariancji warunkowej, znanej względem informacji dostępnych w danym momencie, oraz wariancji bezwarunkowej, uśrednionej po długim okresie. W modelu GARCH zmienność może być przewidywalna warunkowo, mimo że innowacje standaryzowane zachowują się jak biały szum. W standardowym GARCH(1,1) przy warunku skończonej wariancji prognoza wariancji na wiele kroków stopniowo zbliża się do poziomu długookresowego; po nowym dużym szoku prognoza krótkoterminowa rośnie. To wnioskowanie jest warunkowe wobec oszacowanych parametrów i założonej postaci modelu. Niepewność parametrów, zmiana sposobu handlu i ograniczenia jakości danych mogą pogarszać prognozę. W zastosowaniu naukowym warto zatem raportować diagnostykę reszt, błędy prognoz poza próbą i alternatywne specyfikacje, zamiast traktować pojedyncze oszacowanie jako dokładny pomiar ryzyka.

\subsection{8. Skalowanie, fraktale i multifraktalna struktura szeregów finansowych}

Skalowanie oznacza, że pewne wielkości statystyczne zmieniają się z horyzontem czasu według prawa potęgowego w pewnym zakresie skal. Dla przyrostów logarytmicznej ceny \(X(t+\tau)-X(t)\) bada się na przykład momenty bezwzględne
\[
S_q(\tau)=\mathbb E\!\left[|X(t+\tau)-X(t)|^q\right]
\propto \tau^{\zeta(q)}.
\]
Dla samopodobnego procesu z jednym wykładnikiem Hursta \(H\) zachodzi \(\zeta(q)=qH\), jeśli odpowiednie momenty istnieją i przybliżenie potęgowe obowiązuje. Ruch Browna ma \(H=1/2\) i \(\zeta(q)=q/2\). Wartość \(H>1/2\) może wskazywać persystencję, a \(H<1/2\) antypersystencję w analizowanym zakresie; nie jest to samo co korelacja prostych zwrotów w każdym horyzoncie.

Fraktalny opis dotyczy niezmienniczości statystycznej względem zmiany skali i może być stosowany do trajektorii lub miar czasowych. W finansach przykładem badanego efektu jest skalowanie rozkładów zwrotów: rozkłady po przeskalowaniu przez \(\tau^H\) mogą w ograniczonym przedziale horyzontów częściowo się pokrywać. Nie należy jednak utożsamiać takiej obserwacji z dowodem, że rynek jest dokładnie samopodobny we wszystkich skalach. Granice zakresu skal, częstotliwość notowań, sezonowość, nieciągłości i skończona długość próby mają zasadnicze znaczenie.

Multifraktalność oznacza nieliniową zależność \(\zeta(q)\) od \(q\): różne rzędy momentów ujawniają różne wykładniki, co wskazuje na heterogeniczność skalowania i wielorakie źródła zmienności. Można ją badać przez funkcje struktury, analizę multifraktalną detrendowaną (MF-DFA) lub spektrum singularności \(f(\alpha)\), przy czym metody mają różne założenia i podatność na trendy. Momentów wysokiego rzędu nie zawsze da się wiarygodnie oszacować, bo dominują w nich nieliczne ekstremalne zwroty. Pozorną multifraktalność mogą tworzyć grube ogony, niestacjonarność, zmiany reżimu, mikrostruktura albo skończoność próby. Dlatego należy stosować testy na danych zastępczych, analizować zakres skal i podawać niepewność wykładników. Wykładniki skalowania są opisem empirycznym, nie gwarantują prognoz ani nie ustanawiają uniwersalnej teorii rynku.

\end{document}